\documentclass[11pt,reqno]{amsart}
\usepackage{amsmath, amssymb, amsthm}
\usepackage{tensor}
\usepackage{enumerate}
\usepackage[colorlinks=true,allcolors=blue]{hyperref}
\usepackage{blkarray}
\usepackage[numbers]{natbib}
\usepackage[font=footnotesize]{caption}

\usepackage{tikz}
\usetikzlibrary{cd,arrows,calc}

\newcommand{\C}{\mathbb{C}}
\newcommand{\R}{\mathbb{R}}
\newcommand{\Z}{\mathbb{Z}}
\newcommand{\N}{\mathbb{N}}
\newcommand{\Q}{\mathbb{Q}}

\newcommand{\bK}{\mathbb{K}}
\newcommand{\bT}{\mathbb{T}}
\newcommand{\g}{\mathsf{g}}
\newcommand{\liet}{\mathfrak{t}}
\newcommand{\lieg}{\mathfrak{g}}
\newcommand{\affW}{\widetilde{W}}
\newcommand{\extW}{\widehat{W}}

\newcommand{\cR}{\mathcal{R}}
\newcommand{\cRQ}{\mathcal{R}_{\mathbb{Q}}}

\newcommand{\Cuntz}[1]{\mathcal{O}_{#1}}
\newcommand{\Cliff}[1]{\C\ell(#1)}
\newcommand{\Cln}[1]{\C\ell_{#1}}
\newcommand{\Fock}[1]{\mathcal{F}(#1)}

\newcommand{\Aut}[1]{\text{\normalfont Aut}(#1)}
\newcommand{\AutId}[1]{\operatorname{Aut}_0(#1)}
\newcommand{\AutSt}[2]{\operatorname{Aut}_{#1}(#2)}
\newcommand{\eqAut}[2]{\operatorname{Aut}_{#1}(#2)}
\newcommand{\eqAutId}[2]{\operatorname{Aut}_{#1,0}(#2)}
\newcommand{\Autf}[3]{\operatorname{Aut}_{#1}^{#2}(#3)}
\newcommand{\eqIsom}[3]{\operatorname{Iso}_{#1}\left(#2, #3\right)}
\newcommand{\Endo}[1]{\operatorname{End}\left(#1\right)}
\newcommand{\Fun}[1]{\operatorname{Fun}(#1)}

\newcommand{\Proj}[2]{\operatorname{Proj}_{#1}(#2)}
\newcommand{\Rep}[1]{\operatorname{Rep}(#1)}

\newcommand{\id}[1]{\operatorname{id}_{#1}}
\newcommand{\obj}[1]{\text{obj}(#1)}
\newcommand{\EV}[1]{\text{EV}(#1)}
\newcommand{\cA}{\mathcal{A}}
\newcommand{\cP}{\mathcal{P}}
\newcommand{\cC}{\mathcal{C}}
\newcommand{\cD}{\mathcal{D}}
\newcommand{\cG}{\mathcal{G}}
\newcommand{\cH}{\mathcal{H}}
\newcommand{\cV}{\mathcal{V}}
\newcommand{\cI}{\mathcal{I}}
\newcommand{\cU}{\mathcal{U}}

\newcommand{\cEG}[1]{\mathcal{E}_{#1}}
\newcommand{\cBG}[1]{\mathcal{B}_{#1}}

\newcommand{\sfX}{\mathsf{X}}

\newcommand{\cEO}{\mathcal{E}^{\mathcal{O}}} % Cuntz Fell bundle 
\newcommand{\cF}{\mathcal{F}} % higher twist over torus
\newcommand{\cM}{\mathcal{M}} % Morita equivalence for multiplicative structure

\newcommand{\Eig}[2]{\text{Eig}(#1,#2)}
\newcommand{\op}{\text{op}}
\newcommand{\End}{\text{End}}

\newcommand{\lscal}[3]{\tensor*[_{#1}]{\langle #2, #3 \rangle}{}}
\newcommand{\rscal}[3]{\tensor*{\langle #1, #2 \rangle}{_{#3}}}

\newcommand{\rcpt}[2]{\mathcal{K}_{\circ #1}\left(#2\right)}
\newcommand{\lcpt}[2]{\mathcal{K}_{#1 \circ}\left(#2\right)}
\newcommand{\rbdd}[2]{\mathcal{L}_{\circ #1}\left(#2\right)}
\newcommand{\lbdd}[2]{\mathcal{L}_{#1 \circ}\left(#2\right)}

\newcommand{\bddR}{R_{\text{\normalfont bdd}}}

\newcommand{\bdop}[1]{\mathcal{B}{(#1)}}

\newcommand{\Vfin}{\mathcal{V}^{{\rm fin}}_{\C}}
\newcommand{\Viso}{\mathcal{V}^{{\rm iso}}_{\C}}

\newcommand{\MFa}{\mathsf{M}_{\mathnormal{F}}}
\newcommand{\MF}{\mathsf{M}_\mathnormal{F}^{\infty}}

%\newcommand{\VecCf}{\text{Vec}_{\C}^\text{fin}}
\DeclareMathOperator{\Ad}{Ad}
\DeclareMathOperator{\Tor}{Tor}
\DeclareMathOperator{\colim}{colim}
\DeclareMathOperator{\hocolim}{hocolim}
\DeclareMathOperator{\hofib}{hofib}

\newcommand{\grS}{\widehat{\mathcal{S}}}

\newcommand{\Top}{\mathcal{T}op}
\newcommand{\GTop}{G\mathcal{T}op}
\newcommand{\GCStar}{G\text{-}C^*\textbf{Alg}}

\newtheorem{theorem}{Theorem}[section]
\newtheorem{lemma}[theorem]{Lemma}
\newtheorem{corollary}[theorem]{Corollary}
\newtheorem{prop}[theorem]{Proposition}

\theoremstyle{definition}
\newtheorem{definition}[theorem]{Definition}
\newtheorem{remark}[theorem]{Remark}
\newtheorem{example}[theorem]{Example}

\newcommand*\circled[1]{\tikz[baseline=(char.base)]{
            \node[shape=circle,draw,inner sep=0.8pt] (char) {#1};}}
            
\makeatletter
\newcommand{\extp}{\@ifnextchar^\@extp{\@extp^{\,}}}
\def\@extp^#1{\mathop{\bigwedge\nolimits^{\!#1}}}
\makeatother

\begin{document}
\title[Units of $\Z/p\Z$-equivariant $K$-theory and bundles of UHF-algebras]{Units of $\Z/p\Z$-equivariant $K$-theory \\ and bundles of UHF-algebras}
\author{Valerio Bianchi \and Ulrich Pennig}
\address{Cardiff University, School of Mathematics, Senghennydd Road, Cardiff, CF24 4AG, Wales, UK}
\email{BianchiV1@cardiff.ac.uk}
\email{PennigU@cardiff.ac.uk}

\begin{abstract}
We consider infinite tensor product actions of $G = \Z/p\Z$ on the UHF-al\-ge\-bra $D = \Endo{V}^{\otimes \infty}$ for a finite-dimensional unitary $G$-repre\-sen\-ta\-tion $V$ and determine the equivariant homotopy type of the group $\Aut{D \otimes \bK}$, where~$\bK$ are the compact operators on $\ell^2(G) \otimes H_0$ for a separable Hilbert space $H_0$ with $\dim(H_0) = \infty$. We show that this group carries an equivariant infinite loop space structure revealing it as the first space of a naive $G$-spectrum, which we prove to be equivalent to the positive units $gl_1(KU^D)_+$ of equivariant $KU^D$-theory. Here, $KU^D$ is a $G$-spectrum representing $X \mapsto K_*^G(C(X) \otimes D)$. As a consequence the first group of the cohomology theory associated to $gl_1(KU^D)_+$ classifies equivariant $D \otimes \bK$-bundles over finite CW-complexes. 
\end{abstract}

\maketitle

\section{Introduction}

Let $R$ be a commutative ring. Forgetting its additive structure we may focus on the multiplication in $R$ and consider its group of units $GL_1(R)$. The analogue of commutative rings in stable homotopy theory are commutative ring spectra. Famous examples here include the ones representing ordinary cohomology or topological $K$-theory. Similarly to the algebraic setting an associative ring spectrum $E$ has a space of units $GL_1(E)$, which turns out to be the zeroth space of a spectrum $gl_1(E)$ of units in case $E$ is an $E_\infty$-ring spectrum. This theory of units originated in Sullivan's MIT notes \cite{paper:Sullivan}, where they were used to study obstructions to the orientability of vector bundles and spherical fibrations. 

Each model of ring spectra has its own definition of the corresponding units. The definition for $E_\infty$-ring spectra goes back to May, Quinn, Ray and Tornehave \cite{book:May}. The units of a commutative symmetric ring spectrum were initially defined in work by Schlichtkrull \cite{paper:Schlichtkrull}. It was shown by Lind in~\cite{paper:Lind} that the comparison functors between the various models of ring spectra give rise to weak equivalences between their respective units. 

The motivation for the current paper was an observation that interlinks the unit spectrum $gl_1(KU^\cD)$ of a variant of complex topological $K$-theory  with the theory of $C^*$-algebras. Here, $\cD$ denotes a strongly self-absorbing $C^*$-algebra. A unital $C^*$-algebra $\cD$ belongs to this class if there is an isomorphism $\varphi \colon \cD \to \cD \otimes \cD$ that is approximately unitarily equivalent to the left tensor embedding (see Definition~\ref{ssa}). These algebras play a cornerstone role in the classification programme of separable nuclear simple $C^*$-algebras. Examples include infinite tensor products of matrix algebras (i.e.\ infinite UHF-algebras), the Jiang-Su algebra $\mathcal{Z}$ and the infinite Cuntz algebra $\Cuntz{\infty}$. To each such algebra $\cD$ one can associate a commutative symmetric ring spectrum $KU^\cD_*$ that represents the cohomology theory 
\[
    X \mapsto K_*(C(X) \otimes \cD) 
\]
on finite CW-complexes \cite[Sec.~4.1]{paper:DP-Units}, i.e.\ the operator-algebraic $K$-theory of the $C^*$-algebra of continuous $\cD$-valued functions on $X$. In particular, we have weak equivalences $KU^\C \simeq KU^\mathcal{Z} \simeq KU^{\Cuntz{\infty}}$ induced by the unit maps $\C \to \mathcal{Z}$ and $\C \to \Cuntz{\infty}$, and these three spectra represent complex topological $K$-theory $X \mapsto K^*(X)$ by the Serre-Swan theorem. 

Let $\bK$ denote the $C^*$-algebra of compact operators on a separable infinite-dimensional Hilbert space. Together with Dadarlat the second author constructed a map $\Aut{\cD \otimes \bK} \to \Omega KU^\cD_1$ in \cite{paper:DP-Units}, which factors through 
\begin{equation} \label{eqn:noneq_Aut_GL1_Map}
    \Aut{\cD \otimes \bK} \to GL_1(KU^\cD)
\end{equation} 
and extends to a map of infinite loop spaces. The abelian group $K_0(\cD)$ has a natural order structure and the spectrum $gl_1(KU^\cD)$ restricts to the spectrum of positive units $gl_1(KU^\cD)_+$ by pullback via the group homomorphism $GL_1(K_0(\cD))_+ \to GL_1(K_0(\cD))$. The map \eqref{eqn:noneq_Aut_GL1_Map} gives a weak equivalence of infinite loop spaces $\Aut{\cD \otimes \bK} \simeq \Omega^\infty gl_1(KU^\cD)_+$. In case $\cD$ is purely infinite the order structure on $K_0(\cD)$ is trivial (i.e.\ $K_0(\cD)_+ = K_0(\cD)$) and we have $\Aut{\cD \otimes \bK} \simeq \Omega^\infty gl_1(KU^\cD)$ \cite[Thm.~1.1]{paper:DP-Units}. 

Important consequences of this result arise for twisted $K$-theory, which was initially developed in \cite{paper:DonovanKaroubi}, see \cite{paper:KaroubiSurvey} for a survey. The infinite loop map from the last paragraph gives rise to an equivalence
\[
    B\Aut{\cD \otimes \bK} \simeq BGL_1(KU^\cD)_+\ .
\] 
Since the right hand side is the classifying space of (positive) twists of~$KU^\cD$, we obtain a description of twisted $KU^\cD$-theory in terms of bundles of strongly self-absorbing $C^*$-algebras. It includes the geometric twists classified by $K(\Z,3) \simeq B\Aut{\bK}$, i.e.\ the group $[X, K(\Z,3)] \cong H^3(X,\Z)$, considered in \cite{paper:AtiyahSegal, paper:Rosenberg} as a special case. 

Equivariant twisted $K$-theory (with an equivariant geometric twist) has found applications in mathematical physics through a deep theorem by Freed, Hopkins and Teleman \cite{paper:FHT3}: For a compact, simple and simply-con\-nec\-ted Lie group $G$ and $k \in \Z$ the isomorphism classes of positive energy representations of the loop group~$LG$ generate the Verlinde ring $\text{Ver}_k(G)$, which features in conformal and topological field theories. By \cite[Theorem~1]{paper:FHT3} there is a ring isomorphism 
\[
    \text{Ver}_k(G) \cong K^{\dim(G), \tau(k)}_G(G)\ ,
\]
with the fusion product on the Verlinde ring and the Pontrjagin product on twisted $K$-theory. The equivariant twist $\tau(k) \in H^3_G(G,\Z) \cong \Z$ corresponds to $k + h^\vee$, where $h^\vee$ denotes the dual Coxeter number of $G$. 

While it is not too difficult to see that the equivariant geometric twists correspond to cocycles representing classes in $H^3_G(G,\Z)$, a classification of the most general equivariant twists requires the construction of a $G$-space $BGL_1(KU)$ and has not been established. At the point of writing the only results in the literature covering units of equivariant ring spectra are to the best of our knowledge \cite{paper:Santhanam} and \cite[Ex.~5.1.17]{book:SchwedeGlobal}. Nevertheless, candidates for equivariant higher twists over $SU(n)$ in terms of $C^*$-algebra bundles have been constructed in \cite{paper:EvansPennigTwists}.

Apart from the applications in mathematical physics highlighted above, an operator-algebraic model for the units of equivariant $K$-theory will also provide an equivariant refinement of the higher Dixmier-Douady theory in~\cite{paper:DP-Units} with applications in the classification of group actions on $C^*$-algebras: to see why, let $H$ be a countable discrete group and suppose that it acts on $\cD \otimes \bK$ via  
\(
    \alpha \colon H \to \Aut{\cD \otimes \bK}
\).
Now assume that we also have a finite group $G$ that acts through automorphisms on $H$ via $\hat{\gamma}$ and in addition on $\cD \otimes \bK$ in such a way that $\alpha_{\hat{\gamma}_g(h)} = \gamma_g \alpha_h \gamma_g^{-1}$. If we have chosen functorial models for $EH$ and $BH$, then the associated bundle 
\[
    EH \times_{\alpha} \Aut{\cD \otimes \bK} \to BH
\]
comes equipped with a natural $G$-action. Denoting the set of all these $H$-actions with compatible $G$-actions by $\text{Act}_G(H, \cD \otimes \bK)$ we obtain a map
\begin{equation} \label{eqn:action_bundles}
   \text{Act}_G(H, \cD \otimes \bK) \to [BH, B\Aut{\cD\otimes \bK}]^G\ .
\end{equation}
If $\cD$ is a Kirchberg algebra in the UCT-class, $G = \{e\}$ and $H$ is amenable and torsion-free, then the above map induces a bijection between cocycle-conjugacy classes of outer actions and the homotopy set on the right hand side by deep results of Meyer \cite{paper:MeyerGroups}, Gabe and Szab\'{o} \cite{paper:GabeSzabo}. If $BG$ is in addition a finite CW-complex, then this homotopy set is the first group of a cohomology theory making it accessible to computations via algebraic topology.
In the equivariant case the right hand side should evaluate to the equivariant cohomology theory associated to the positive unit spectrum of equivariant $KU^{\cD}$-theory. In fact, \eqref{eqn:action_bundles} makes sense with much weaker assumptions on~$G$. However, equivariant stable homotopy theory becomes more intricate for more general groups.

Developing these ideas requires combining methods from equivariant stable homotopy theory with topological considerations about automorphism groups of $C^*$-algebras. In this paper we present an important first advance in this direction. As a starting point we picked a setting that is very explicit: We fix the group $G$ to be $\Z/p\Z$ for a prime $p$, and we consider a UHF-algebra~$D$ given by the infinite tensor product of a unitary finite-dimensional $G$-representation $V$, i.e.
\[
    D = \Endo{V}^{\otimes \infty}\ .
\]
This is a $G$-$C^*$-algebra in a natural way when equipped with the infinite tensor product of the action by conjugation. The restriction to prime order cyclic groups has two benefits: There are only two subgroups to consider, and some of the results obtained in \cite{paper:EP-Circle} for circle actions on UHF-algebras carry over immediately to the case of $\Z/p\Z$-actions. We then determine the $G$-equivariant homotopy type of $\Aut{D \otimes \bK}$. In particular, we show that the automorphism group is an equivariant infinite loop space associated to a naive $G$-spectrum, which we identify in our main result.
\begin{theorem}[Thm.~\ref{thm:main_thm}]
    Let $G = \Z/p\Z$ for a prime $p \in \N$. Let $V$ be a finite-dimensional unitary $G$-representation and let
    \[
        D = \Endo{V}^{\otimes \infty}
    \]
    be the associated $G$-$C^*$-algebra. The group $\Aut{D \otimes \bK}$ is a $G$-equivariant infinite loop space underlying a naive $G$-spectrum $EH_{\cI}\Aut{D \otimes \bK}$ and we have an equivalence of naive $G$-spectra
    \[
        EH_{\cI}\Aut{D \otimes \bK} \simeq gl_1(KU^D)_+\ .
    \]
\end{theorem}

The group $\Aut{D \otimes \bK}$ has two equivariant deloopings: one arising from the group structure, the other one from the tensor product. Just as in the non-equivariant setting these two agree. As a result the first group of the cohomology theory associated to $gl_1(KU^D)_+$ classifies equivariant $C^*$-algebra bundles with fibre $D \otimes \bK$ in the following sense:
\begin{corollary}[Cor.~\ref{cor:bdl_classification}]
Let $E_D^*(X) = gl_1(KU^D)_+^*(X)$ be the $\Z/p\Z$-equivariant cohomology theory associated to the spectrum $gl_1(KU^D)_+$. Let~$X$ be a finite CW-complex. Then 
\[
    E_D^0(X) \cong [X, \Aut{D \otimes \bK}]^{\Z/p\Z} \quad \text{and} \quad E_D^1(X) \cong [X, B\Aut{D \otimes \bK}]^{\Z/p\Z}\ .
\]
In particular, isomorphism classes of $\Z/p\Z$-equivariant $C^*$-algebra bundles with fibres isomorphic to the $\Z/p\Z$-algebra $D \otimes \bK$ over $X$ form a group with respect to the fibrewise tensor product, which is isomorphic to $E_D^1(X)$. 
\end{corollary}

The article is structured as follows. In Section \ref{prel_1} we describe an infinite loop space machine that is equivariant for finite groups based on diagram spaces, called commutative $\cI$-$G$-monoids, whose output is a naive positive $\Omega$-$G$-spectrum. In Section \ref{prel_2} we turn to operator algebras: we provide some background on strongly self-absorbing $C^*$-algebras, describe our basic setup and recall how Bott periodicity can be phrased in terms of Hilbert module bundles. \\
In Section \ref{sec3} we construct an equivariant spectrum of units for a localisation of equivariant $K$-theory. This is done in two steps: in Section~\ref{spectrum} we refine the construction of a commutative symmetric ring spectrum $KU^D$ associated to a strongly self-absorbing $C^*$-algebra $D$ to the equivariant case (Definition~\ref{def:KUD}). Then we construct its equivariant units in Section~\ref{equnits} at the level of diagram spaces, and finally lift them to a $G$-spectrum of units (Definition~\ref{def:spec_units}). \\
Our goal is to compare this spectrum to the one associated to the automorphism group of the stabilisation of $D=\Endo{V}^{\otimes \infty}$ for an action of $G=\Z/p\Z$. In order to be able to do that, we interpret $\Aut{D \otimes \bK}$ as the first space in a commutative $\cI$-$G$-monoid $\cG_D$ in Section \ref{G_D}. In Section~\ref{Aut} we then collect a series of results which provide the equivariant homotopy type of this group, mainly using techniques from \cite{paper:EP-Circle}. We also compute the homotopy groups of the equivariant automorphisms in Corollary~\ref{cor:pikAut}.\\
The main result of this paper is contained in Section \ref{final_section}, where we explicitly describe a map between the $\Z/p\Z$-spectrum coming from the equivariant infinite loop space structure of $\Aut{D \otimes \bK}$ and the spectrum of $\Z/p\Z$-equivariant units of $KU^D$, and prove that it is an equivalence onto the positive units (Theorem~\ref{thm:main_thm}). In particular, this provides a classification of $\Z/p\Z$-equivariant $\Aut{D \otimes \bK}$-bundles in terms of an equivariant cohomology theory (Corollary~\ref{cor:bdl_classification}). We discuss a model for the classifying space of such equivariantly locally trivial bundles in Section \ref{final_1}, and we state and prove the main theorems in Section \ref{final_2}.\\
Constructing the equivariant classifying space $B\Aut{D \otimes \bK}$ as the geometric realisation of the nerve is only possible, if we can show that $\Aut{D \otimes \bK}$ is equivariantly well-pointed, which is the content of Lemma~\ref{lem:G-cofib} in the appendix. There we also prove a few other topological results that might be of independent interest: for example, we show that $\Aut{D \otimes \bK}$ has the $G$-homotopy type of a CW-complex (Lemma~\ref{lem:GCW-complex}) and that an evaluation map $\Aut{A} \to \Proj{}{A}$ for a finite group $G$ and a $G$-$C^*$-algebra $A$ is a Hurewicz $G$-fibration when suitably restricted in the codomain (Lemma~\ref{lem:Hurewicz}). 

\subsection*{Acknowledgements} 
The authors would like to thank John A.\ Lind, Steffen Sagave and Stefan Schwede for helpful discussions. Some parts of this work emerged from the PhD thesis of the first author, which was funded by the EPSRC grant EP/W52380X/1 (project reference: 2601068).

\section{Preliminaries}

\subsection{Preliminaries from equivariant homotopy theory}
\label{prel_1}

Let $\GTop$ be the category of based compactly generated weakly Hausdorff spaces equipped with an action of a finite group $G$ which fixes the basepoint, and based equivariant maps. Note that $\GTop$ has a model structure whose weak equivalences are morphisms $f:X \rightarrow Y$ such that $f^H: X^H \rightarrow Y^H$ is a weak equivalence in the underlying category $\Top$, for all subgroups $H \subseteq G$.

The orbits $G/H$ form a full subcategory of $\GTop$ which we denote by~$O_G$, the orbit category. Note that for subgroups $K \subseteq H \subseteq G$, we have a natural map of $G$-sets $G/K \rightarrow G/H$, and a corresponding inclusion map of fixed points $X^H \rightarrow X^K$ in the opposite direction for a $G$-space $X$. This means we have a functor
\[O_G^{\op} \rightarrow \Top \qquad \text{ given by } \qquad G/H \mapsto X^H\ .\]
\begin{definition}
    An $O_G$-\textit{space} is a functor $O_G^\op \rightarrow \Top$. Denote by \emph{$O_G$-spaces} the category of $O_G$-spaces and continuous natural transformations.
\end{definition}
As described above there is a fixed point functor $\GTop \rightarrow O_G$-spaces defined on objects by
\[
X \mapsto \big((G/H) \mapsto X^H \big)\ .
\]
By a celebrated theorem by Elmendorf \cite{paper:Elmendorf} this functor is an equivalence on homotopy categories. Its homotopy inverse can be constructed explicitly as the geometric realization of the simplicial bar construction:  
\begin{equation}\label{eqn:elmendorf}
\Phi \colon O_G\text{-spaces} \to \GTop \quad, \quad X \mapsto |B(X,O_G,M)|\ ,
\end{equation}
where $M$ is the forgetful functor realising orbits as spaces (see \cite[Sec. 8.8]{book:Kervaire} for details).

Following \cite{paper:Lind} we now define two diagram categories of $G$-spaces. Let $\mathcal{I}$ be the category of finite sets $\textbf{n}=\{1,\dots,n\}$ (including $\mathbf{0} = \emptyset$) and injective functions. It is a symmetric monoidal category via concatenation of finite ordered sets $\textbf{m} \sqcup \textbf{n}$ along with symmetric structure maps $\textbf{m} \sqcup \textbf{n} \rightarrow \textbf{n} \sqcup \textbf{m}$ given by the obvious shuffles.

\begin{definition}
An $\mathcal{I}$-$G$-\textit{space} is a functor $X \colon \mathcal{I} \rightarrow \GTop$.
An $\mathcal{I}$-$G$-space~$X$ is called a (commutative) $\mathcal{I}$-$G$-\textit{monoid} if it comes equipped with a natural transformation $\mu: X \times X \rightarrow X \circ \sqcup$ which is associative, unital (and commutative) in the sense of \cite[Def.~2.2]{paper:DP-Units}.
\end{definition}

Let $\Gamma^{\op}$ be the category of finite based sets $\textbf{n}^+=\{0,1,\dots,n\}$ with basepoint~$0$ and based maps.

\begin{definition}
A $\Gamma$-$G$-\textit{space} is a functor $\Gamma^{\op} \rightarrow \GTop$. 
A $\Gamma$-$G$-space $A$ is called \textit{special} if for every $n$ the based $G$-map 
\[A(\mathbf{n}^+) \rightarrow \big(\prod_{i=1}^{n}A(\textbf{1}^+)\big)\]
induced by projections is a weak equivalence. We refer to this map as the equivariant Segal map. 
\end{definition}

The notions of $\cI$-$O_G$-\textit{spaces} and (special) $\Gamma$-$O_G$-\textit{spaces} are defined analogously by replacing the target category by $O_G$-spaces.

There is a way to associate a special $\Gamma$-$G$-space to any commutative $\mathcal{I}$-$G$-monoid. This works just as in the non-equivariant case \cite[Construction~12.1]{paper:Lind} with some minor adjustments; for the sake of completeness, we recall how it works.

Let $X$ be a commutative $\cI$-$G$-monoid. Let $P(\textbf{n}^+)$ be the category of finite sets $0 \notin A \subset \textbf{n}^+$ and define $\mathcal{I}(\textbf{n}^+)$ to be the category of functors
$\theta: P(\textbf{n}^+) \rightarrow \mathcal{I}$
with the property that for every pair of disjoint subsets $A,B \in P(\textbf{n}^+)$ the diagram
\[\theta(A) \rightarrow \theta(A \cup B) \leftarrow \theta(B)\]
(induced by inclusions) is a coproduct diagram in $\mathcal{I}$. Morphisms in $\mathcal{I}(\textbf{n}^+)$ are natural transformations.
A functor $\theta \in \mathcal{I}(\textbf{n}^+)$ is given by the data of a collection of objects $\theta_{i}:=\theta(\{i\}) \in \mathcal{I}$ for each $1 \le i \le n$ (by definition, $\theta(\emptyset)=\textbf{0}$), plus morphisms $\theta(i) \rightarrow \theta(A)$ for each $A \subset \{1,\dots,n\}$ which assemble into a canonical isomorphism 
\begin{equation} \label{eqn:theta}
    \bigoplus_{i \in A} \theta_{i} \cong \theta(A)\ .
\end{equation}
Consider the forgetful functor $f: \mathcal{I}(\textbf{n}^+) \rightarrow \mathcal{I}^{n}$ that sends $\theta$ to $(\theta_{1}, \dots, \theta_{n})$. We get a functor $X(\textbf{n}^+):\mathcal{I}(\textbf{n}^+) \rightarrow \GTop$ given by the composite
\[X(\textbf{n}^+): \mathcal{I}(\textbf{n}^+) \xrightarrow{f} \mathcal{I}^{n} \xrightarrow{X^{n}} \GTop\]
\[\theta \mapsto \prod_{i=1}^{n}X(\theta_{i})\ .\]
Finally, define a functor $H_{\mathcal{I}}X: \Gamma^{\op} \rightarrow \GTop$ as the homotopy colimit (in the sense of \cite[Def. 5.8.1]{book:Kervaire})
\[H_{\mathcal{I}}X(\textbf{n}^+)=\hocolim_{\mathcal{I}(\textbf{n}^+)}X(\textbf{n}^+)\]
and set $H_{\mathcal{I}}X(\textbf{0}^+)=*$. Note that $H_{\mathcal{I}}X(\textbf{1}^+)=\hocolim_{\mathcal{I}}X=:X_{h \cI}$.\\
In order to make $H_{\cI}X$ into a $\Gamma$-$G$-space we need to have functoriality; this is only possible because $X$ has the structure of an commutative $\cI$-$G$-monoid. Let $\alpha: \textbf{m}^+ \rightarrow \textbf{n}^+$ be a map of finite based sets. There is a functor $\alpha_{*}: \mathcal{I}(\textbf{m}^+) \rightarrow \mathcal{I}(\textbf{n}^+)$ given by precomposition with $\alpha^{-1}$. We have a natural transformation $X(\alpha): X(\textbf{m}^+) \rightarrow X(\textbf{n}^+) \circ \alpha_{*}$ defined as
\[
\begin{tikzcd}[column sep=0.5cm,row sep=0cm]
X(\alpha)_{\theta} \colon \prod_{i=1}^{m}X(\theta_{i}) \ar[r,"\pi"] & \prod_{j=1}^{n} \prod_{i \in \alpha^{-1}(j)} X(\theta_{i}) \ar[r,"\mu"] & \prod_{j=1}^{n} X \bigl( \bigoplus_{i \in \alpha^{-1}(j)} \theta_{i} \bigr) \\
&&\cong \prod_{j=1}^{n} X(\theta(\alpha^{-1}(j)))
\end{tikzcd}
\]
where $\pi$ is the projection away from the factors indexed by elements $i \in \textbf{m}^+$ that are mapped to the basepoint by $\alpha$, and the last isomorphism is induced by the canonical isomorphism \eqref{eqn:theta}. In this definition a choice is apparently involved when ordering the product indexed by $i \in \alpha^{-1}(j)$, but commutativity of $\mu$ ensures that it does not actually matter. Finally, we lift the natural transformation $X(\alpha)$ to the first map of homotopy colimits below:
\[ \hocolim_{\mathcal{I}(\textbf{m}^+)}X(\textbf{m}^+) \rightarrow \hocolim_{\mathcal{I}(\textbf{m}^+)}X(\textbf{n}^+) \circ \alpha_{*} \xrightarrow{\alpha_{*}} \hocolim_{\mathcal{I}(\textbf{n}^+)}X(\textbf{n}^+)\ .\]
This composition defines $H_{\mathcal{I}}X(\alpha).$

\begin{prop}
The functor $H_{\mathcal{I}}X$ is a special $\Gamma$-$G$-space.
\end{prop}

\begin{proof}
We need to check that the equivariant Segal map
\[H_{\mathcal{I}}X(\mathbf{n}^+) \rightarrow \big(\prod_{i=1}^{n}H_{\mathcal{I}}X(\textbf{1}^+)\big)\]
is a weak equivalence. Using the canonical isomorphism
\[(\hocolim_{\mathcal{D}}X)^{n} \cong \hocolim_{\mathcal{D}^{n}}X^{n}\]
this boils down to checking that the map of homotopy colimits induced by the forgetful functor $f$
\[ \hocolim_{\mathcal{I}(\textbf{n}^+)}X(\textbf{n}^+) \rightarrow \hocolim_{\mathcal{I}^{n}}X^{n}\]
is a weak equivalence, which in turn follows from the Bousfield-Kan cofinality criterion \cite[Thm. 5.8.15]{book:Kervaire} applied to the forgetful functor $f: \mathcal{I}(\textbf{n}^+) \rightarrow \mathcal{I}^{n}$, provided we can show that $f$ is homotopy final in the sense of \cite[Def.~5.8.13]{book:Kervaire}. This reduces to verify that for any $v=(v_{1},\dots,v_{n}) \in \mathcal{I}^{n}$, the comma category $(v/f)$ has an initial object. This is in fact the case, and such an initial object is given by the triple 
\[\bigl(v, \theta, \text{id}: v \rightarrow f(\theta)\bigr) \in (v/f)\ ,\]
where $\theta$ is the functor $P(\textbf{n}^+) \rightarrow \mathcal{I}$ given by
\[
\theta(A)=\bigoplus_{i \in A} v_i\ .\qedhere
\]
\end{proof}

\begin{remark}
    In the rest of the paper we will often work with $O_G$-spa\-ces instead of $G$-spaces. The construction outlined above can be repeated verbatim when $X$ is a commutative $\cI$-$O_G$-monoid. The output in this case is a special $\Gamma$-$O_G$-space, which we will always turn into a special $\Gamma$-$G$-space again using $\Phi$ from \eqref{eqn:elmendorf}. Hence, we will consider $H_{\cI}$ as a functor from commutative $\cI$-$O_G$-monoids to $\Gamma$-$G$-spaces. Here a key observation is that the functor $\GTop \rightarrow O_G$-spaces preserves products on the nose, hence allowing us to compare the output of this machinery to the one that we get by starting with the corresponding commutative $\cI$-$G$-monoid. As one would expect the two resulting objects are homotopy equivalent, because of Elmendorf's theorem and the fact that the homotopy colimit construction commutes with taking fixed points.
\end{remark}
    
There are several ways to get a $G$-spectrum out of a special $\Gamma$-$G$-space. All of these constructions, known in the literature as infinite loop space machines, require some sort of $G$-cofibrancy condition on the $\Gamma$-$G$-space and differ by the flavour of the $G$-spectra they produce. 
For our purposes it is enough to get a naive $G$-spectrum; in this case, the cofibrancy condition can always be assumed to hold thanks to a functorial cofibrant replacement theorem 
\cite[Prop. 2.12]{paper:MayMerlingOsorno}. The following is a straightforward generalization of the Segal machine to $G$-spaces \cite[Def. 2.22]{paper:MayMerlingOsorno}.

\begin{definition}
    Let $A$ be a special $\Gamma$-$G$-space. For each $p$, let $A[\textbf{p}^+]$ be the $\Gamma$-$G$-space that sends $\textbf{q}^+$ to $A(\textbf{p}^+ \wedge \textbf{q}^+)$. Define the classifying $\Gamma$-$G$-space $\mathbb{B}A$ to be the $\Gamma$-$G$-space whose $p$-th $G$-space is the geometric realization $|A[\textbf{p}^+]|$. Now let $\mathbb{B}^0A=A$ and iteratively define $\mathbb{B}^nA=\mathbb{B}(\mathbb{B}^{n-1}A)$. We denote by $EA$ the resulting naive $G$-spectrum with $n$-th $G$-space 
    \[(EA)_n=(\mathbb{B}^nA)(\mathbf{1}^+)\ .\]
\end{definition}

Note that this machine actually produces positive $\Omega$-$G$-spectra, and the first structure map $(EA)_0=A(\textbf{1}^+) \rightarrow \Omega (EA)_1$ is a group completion \cite[Prop. 2.18]{paper:MayMerlingOsorno}.

So far we have shown how to construct a positive $\Omega$-$G$-spectrum out of a commutative $\cI$-$O_G$-monoid. Our goal for the rest of the subsection is to show that for a stable $\cI$-$O_G$-monoid $X$ (see Definition \ref{stable}), all of the information about the homotopy type of such a spectrum (except in level $0$) is contained in the $O_G$-space $X(\textbf{1})$.

The naive $G$-spectra we obtain as the output of the equivariant Segal machine (using $\Gamma$-$G$-spaces) are spectra objects in $G$-spaces. Therefore homotopy groups and equivalences are defined as follows:
\begin{definition}
Let $n$ be an integer and $H \subseteq G$ be a subgroup. The $n$-th equivariant homotopy group $\pi_n^H X$ of a naive $G$-spectrum $X$ is defined as
\[\pi_n^H(X) = \colim_{k \to \infty} \left[S^{n + k}, X_k\right]^H \]
where $H$ acts trivially on $S^{n+k}$. We say a morphism $f:X \rightarrow Y$ of naive $G$-spectra is an equivalence if the induced map $\pi_n^H(f):\pi_n^H(X) \rightarrow \pi_n^H(Y)$ is an isomorphism for all integers $n$ and all subgroups $H \subseteq G$.

    %Fix a $G$-set universe $\cU$ and let $n$ be an integer and $H \subseteq G$ be a subgroup. The $n$-th equivariant homotopy group $\pi_n^H X$ of a symmetric $G$-spectrum $X$ is defined as
    %\[\pi_n^H X= \colim_{M \in s_G(\cU)} [S^{n \sqcup M}, X(M)]^H\ ,\]
    %where $s_G(\cU)$ is the poset of finite $G$-subsets of $\cU$ (partially ordered by inclusion).
\end{definition}

\begin{lemma}
\label{+omega}
For every positive $\Omega$-$G$-spectrum $X$, every $n \ge 0$ and every subgroup $H \subseteq G$ the map
\[\pi_n(\Omega X_1^H) \rightarrow \pi_n^H(X)\]
is an isomorphism.
\end{lemma}

\begin{proof}
Note that the right hand side is $\pi_n^H(\Omega^k X_k)$ and by the definition of a positive $\Omega$-$G$-spectrum we have that $\Omega X_1 \to \Omega^{k} X_{k}$ is an equivariant weak equivalence for $k \ge 1$.
%See \cite[Ex. 3.2]{paper:Hausmann} with $X(M)=X_1$.
\end{proof}

\begin{definition}
\label{stable}
    An $\cI$-$O_G$-space is called stable if all morphisms $\mathbf{m} \to \mathbf{n}$ in~$\cI$ with $m > 0$ are mapped to weak equivalences for all cosets~$G/H$.
\end{definition}

\begin{remark}
\label{rem_stable}
    The category $\cI$ is filtered and the homotopy group functor $\pi_*$ preserves filtered colimits. Moreover, taking fixed points commutes with homotopy colimits for actions of finite groups. This implies that for a stable $\cI$-$O_G$-space $Z$ we have
    \[
    Z(\textbf{1})(G/H) \simeq \Phi(Z(\mathbf{1}))^H \simeq ((\Phi\circ Z)_{h\cI})^H \simeq Z(G/H)_{h\cI}\ .\]
\end{remark}

\begin{lemma}
\label{deloop}
Let $X, Y$ be two \textit{stable} commutative $\cI$-$O_G$-monoids such that $\pi_0(X(\mathbf{1})(G/H))$ and $\pi_0(Y(\mathbf{1})(G/H))$ are groups for all subgroups $H \subseteq G$ and let $f: X \rightarrow Y$.
We have the following commutative diagram
% https://q.uiver.app/#q=WzAsNixbMCwwLCJcXHBpX3trfV57SH0oRUhfe1xcbWF0aGNhbHtJfX1YKSJdLFsyLDAsIlxccGlfe2t9XntIfShFSF97XFxtYXRoY2Fse0l9fVkpIl0sWzAsMSwiXFxwaV97a30gKChFSF97XFxtYXRoY2Fse0l9fVhfezF9KV57SH0pIl0sWzIsMSwiXFxwaV97a30gKChFSF97XFxtYXRoY2Fse0l9fVlfezF9KV57SH0pIl0sWzAsMiwiXFxwaV97a30oWCgxKV57SH0pIl0sWzIsMiwiXFxwaV97a30oWSgxKV57SH0pIl0sWzAsMV0sWzIsMCwiXFxzaW1lcSJdLFsyLDNdLFszLDEsIlxcc2ltZXEiLDJdLFs0LDVdLFs1LDMsIlxcc2ltZXEiLDJdLFs0LDIsIlxcc2ltZXEiXV0=
\[\begin{tikzcd}
	{\pi_{n}^{H}(EH_{\mathcal{I}}X)} && {\pi_{n}^{H}(EH_{\mathcal{I}}Y)} \\
	{\pi_{n+1} (((EH_{\mathcal{I}}X)_1)^{H})} && {\pi_{n+1} (((EH_{\mathcal{I}}Y)_1)^{H})} \\
	{\pi_{n}(X(\mathbf{1})(G/H))} && {\pi_{n}(Y(\mathbf{1})(G/H))}
	\arrow[from=1-1, to=1-3]
	\arrow["\simeq", from=2-1, to=1-1]
	\arrow[from=2-1, to=2-3]
	\arrow["\simeq"', from=2-3, to=1-3]
	\arrow[from=3-1, to=3-3]
	\arrow["\simeq"', from=3-3, to=2-3]
	\arrow["\simeq", from=3-1, to=2-1]
\end{tikzcd}\]
In particular, $f$ induces an isomorphism on $\pi_n^H$ of the corresponding $G$-spectra if and only if the map
\[\pi_{n}(X(\mathbf{1})(G/H)) \rightarrow \pi_{n}(Y(\mathbf{1})(G/H))\]
is an isomorphism.
\end{lemma}

\begin{proof}
The two upper vertical arrows are isomorphisms by Lemma \ref{+omega}. Now let $Z \in \{X,Y\}$. Note that $\pi_0^H(Z(\mathbf{1})) \cong \pi_0^H(Z_{h\cI}) \cong \pi_0^H((EH_\cI Z)_0)$ is a group by assumption. Observe that $((EH_\cI Z)_0)^H \simeq (\Omega (EH_{\cI} Z)_1)^H$ by the group completion theorem and 
\[((EH_{\mathcal{I}}Z)_0)^H \simeq (H_{\mathcal{I}}Z(\textbf{1}^+))^H \simeq ((\Phi \circ Z)_{h\cI})^H\ .\]
Thus the two lower vertical arrows are isomorphisms as well by Remark~\ref{rem_stable}.
\end{proof}

\subsection{Preliminaries from operator algebras.} \label{prel_2}

We will look at $G$-$C^*$-alge\-bras that arise from infinite tensor products of $\Z/p\Z$-representations. Such algebras belong to the class of strongly self-absorbing $C^*$-algebras \cite{paper:TomsWinter}.
\begin{definition}
\label{ssa}
A separable, unital $C^*$-algebra $D$ is called \textit{strongly self-absorbing} if it tensorially absorbs itself, via an isomorphism $\psi: D \rightarrow D \otimes D$, in such a way that $\psi$ is asymptotically unitarily equivalent to the left tensor embedding $d \mapsto d \otimes 1$.    
\end{definition}

To put this class into context note that some properties that are essential in the classification programme of nuclear, separable, simple $C^*$-algebras are related to tensorial absorption of strongly self-absorbing ones:
\begin{enumerate}[(i)]
    \item If $A$ is separable, simple and nuclear, then $A$ is purely infinite if and only if $A \otimes \Cuntz{\infty} \cong A$ \cite[Thm.~3.15]{paper:KirchbergPhillips}, \cite[Thm.~7.2.6 (ii)]{book:RordamStormer}. Here, $\Cuntz{\infty}$ denotes the infinite Cuntz algebra, which is strongly self-absorbing (and therefore purely infinite).  
    \item Separable, simple, unital, nuclear $C^*$-algebras in the UCT class with at most one trace are classified up to $\mathcal{Z}$-stability by their ordered $K$-theory \cite[Cor.~E]{paper:TWW}. Here, $\mathcal{Z}$ denotes the Jiang-Su algebra, which is strongly self-absorbing and ``up to $\mathcal{Z}$-stability'' means after taking a tensor product with $\mathcal{Z}$.
    \item Let $A$ be a simple, separable, unital and nuclear $C^*$-algebra. Then $A \otimes \Cuntz{2} \cong \Cuntz{2}$ \cite[Cor.~3.8]{paper:KirchbergPhillips}. Here, $\Cuntz{2}$ is the Cuntz algebra on two generators, which is also strongly self-absorbing.  
\end{enumerate}
Apart from $\C$, $\mathcal{Z}$, $\Cuntz{\infty}$ and $\Cuntz{2}$ other examples of strongly self-absorbing $C^*$-algebras arise from UHF-algebras, which are infinite tensor products of matrix algebras: Let $n \in \N$ and consider the sequence
\[ 
   \C \to \dots \to M_n(\C)^{\otimes k} \to M_n(\C)^{\otimes (k+1)} \to M_n(\C)^{\otimes (k+2)} \to \dots\ ,
\]
where the connecting homomorphisms are given by $T \mapsto T \otimes 1_n$. Denote the colimit of this sequence in the category of unital $C^*$-algebras by $M_n^{\otimes \infty}$. If $p_1, \dots, p_r$ are the prime factors of $n$, then we have 
\[
    M_n^{\otimes \infty} \cong M_{p_1}^{\otimes \infty} \otimes \dots \otimes M_{p_r}^{\otimes \infty}\ .
\]
For an arbitrary subset $P$ of the prime numbers we therefore define (by slight abuse of notation) 
\[
    M_P^{\otimes \infty} = \bigotimes_{p \in P} M_p^{\otimes \infty}\ .
\]
If $P$ is the set of all prime numbers, the resulting algebra is called the universal UHF-algebra, usually denoted by $\mathcal{Q}$, because $K_0(\mathcal{Q}) \cong \Q$. A list of all strongly self-absorbing $C^*$-algebras in the UCT-class can be found in Figure~\ref{fig:ssa_diamond}.
\begin{figure}[htp]
    \centering
    \begin{tikzcd}[column sep=1.4cm,row sep=0.2cm]
        & \mathcal{Z} \ar[dd] \ar[r] & M_P^{\otimes \infty} \ar[r] \ar[dd] & \mathcal{Q} \ar[dd] \ar[dr] \\ 
        \C \ar[ur] \ar[dr] &&&& \Cuntz{2} \\
        & \Cuntz{\infty} \ar[r] & M_P^{\otimes \infty} \otimes \Cuntz{\infty} \ar[r] & \mathcal{Q} \otimes \Cuntz{\infty} \ar[ur]
    \end{tikzcd}
    \caption{\label{fig:ssa_diamond}The list of strongly self-absorbing $C^*$-algebras in the UCT class. An arrow indicates tensorial absorption (e.g.\ $\Cuntz{\infty} \otimes \mathcal{Z} \cong \Cuntz{\infty}$). }
    \label{fig:ssa_diamond}
\end{figure}
%\pagebreak{10}

Let $D$ be a strongly self-absorbing $C^*$-algebra and consider the functor $X \mapsto K_*(C(X) \otimes D)$. The tensor product induces a ring structure on $K_*(C(X) \otimes D)$, given by
\[
\begin{tikzcd}[column sep=0.6cm,every matrix/.append style={nodes={font=\small}}]
    K_*(C(X) \otimes D) \times K_*(C(X) \otimes D) \ar[r] & K_*(C(X \times X) \otimes D) \ar[r,"\Delta^*"] & K_*(C(X) \otimes D)
\end{tikzcd}
\]
where the first arrow takes the tensor product and applies the inverse isomorphism $\psi^{-1} \colon D \otimes D \to D$, and $\Delta \colon X \to X \times X$ denotes the diagonal map. The multiplication does not depend on the choice of $\psi$ because $\Aut{D}$ is contractible \cite[Thm. 2.3]{paper:DP-Dixmier} and $K$-theory is homotopy invariant. Thus, the assignment $X \rightarrow K_*(C(X) \otimes D)$ defines a multiplicative cohomology theory on finite CW-complexes.

Because of its infinite tensor product structure it is easy to construct interesting group actions on UHF-algebras. Let $G$ be a discrete group and let~$\rho \colon G \to U(V)$ be a finite-dimensional unitary $G$-representation, then $G$ acts by conjugation on $\Endo{V}$. This $G$-action extends to the UHF-algebra
\[
    D = \Endo{V}^{\otimes \infty}\ .
\]
Note that while the action of $G$ on $\Endo{V}$ is inner, this is no longer the case for the induced action on $D$, where the action is in general only approximately inner. In the next section we will combine the above two observatios and construct a $G$-spectrum $KU^D$ representing the equivariant cohomology theory $X \mapsto K_*^G(C(X) \otimes D)$.

We will also need the Swan theorem in $K$-theory with coefficients in a $C^*$-algebra $A$: Recall that to any projection valued function in $C(X,M_n(A))$ one can associate a finitely generated projective Hilbert $A$-module bundle, and vice versa (see \cite[Thm. 3.14]{paper:Schick}). If we denote by $K(X,A)$ the Grothendieck group of isomorphism classes of such bundles over a compact Hausdorff space~$X$, we have $K(X,A) \cong K_0(C(X) \otimes A)$ \cite[Prop. 3.17]{paper:Schick}. As an example, we illustrate how this correspondence allows for a more intuitive description of the index map in $K$-theory.

\begin{example}
\label{indexmap}
Let $A$ be a unital $C^*$-algebra. By Bott periodicity $K_1(A)$ is isomorphic to $K_0(SA)$ where $SA=C_0((0,1),A)$ is the suspension of~$A$. This isomorphism has the following explicit description (see for instance \cite[Thm. 10.1.3]{book:RLL}): let $u \in M_{2n}(A)$ and choose a path 
\[
    \gamma \colon (0,1) \rightarrow U(M_{2n}(A)) \quad \text{with } \gamma(0)=
\begin{pmatrix}
1 & 0\\
0 & 1
\end{pmatrix}
\text{ and } \gamma(1)=
\begin{pmatrix}
u & 0\\
0 & u^*
\end{pmatrix}\ . 
\]
The image of $[u] \in K_1(A)$ is the formal difference of projections 
\[
\bigg[\gamma 
\begin{pmatrix}
1 & 0\\
0 & 0
\end{pmatrix}
\gamma^*\bigg] - \bigg[
\begin{pmatrix}
1 & 0\\
0 & 0
\end{pmatrix}\bigg]\ ,
\]
which defines an element in $K_0(SA)$. This element corresponds to the formal difference $[\mathcal{H}_{\gamma}]-[\underline{\mathcal{H}}]$ of Hilbert $A$-module bundles over $S^1$, where
\[\mathcal{H}_{\gamma}=\{([t],v) \in S^1 \times \C^2 \otimes A \,\,| \,\,\gamma(t)^*
\begin{pmatrix}
1 & 0\\
0 & 0
\end{pmatrix}
\gamma(t)v=v\}\]
and $\underline{\mathcal{H}}$ is the trivial bundle. To the same unitary $u \in U(A)$ we can also associate another Hilbert $A$-module bundle over $S^1$, defined by
\[\mathcal{H}_u=(0,1) \times A\,/ \sim \ , \qquad \text{where } (1,a) \sim (0,ua)\ .\]
The upshot is that this simpler construction already encodes the suspension isomorphism (up to a basepoint shift). More precisely, there is an isomorphism of Hilbert $A$-module bundles
\[\mathcal{H}_u \xrightarrow{\cong} \mathcal{H}_{\gamma}\ ,\]
defined by $(t,a) \mapsto ([t],\gamma(t)^*(e_1 \otimes a))$.
\end{example}

\section{The $G$-spectrum of units of $KU^D$-theory}
\label{sec3}

\subsection{The spectrum $KU^D$}
\label{spectrum}
In this section we will construct a $G$-equivariant commutative symmetric ring spectrum $KU^D$ for $G = \Z/p\Z$ associated to a functor taking values in graded $G$-$C^*$-algebras, i.e.\ $\Z/2\Z$-graded $C^*$-algebras together with an action by the group $G$ that preserves the grading. 

\begin{definition}
    Let $\Sigma$ be the category of finite sets and bijective maps. This is a symmetric monoidal category with respect to the coproduct $\sqcup$. Let $\GCStar$ be the category of graded $G$-$C^*$-algebras and equivariant grading-preserving $*$-homomorphisms. This is also a symmetric monoidal category with respect to the graded minimal tensor product $\otimes$. A \textit{coefficient system} is a symmetric monoidal functor 
    \[
        C \colon (\Sigma, \sqcup) \to (\GCStar, \otimes)\ .
    \]
    The tensor product of two coefficient systems $C_1$ and $C_2$ is defined to be  
    \[
        (C_1 \otimes C_2)(S) = C_1(S) \otimes C_2(S)
    \]
    for a finite set $S$. Let $\eta^i_{S,T} \colon C_i(S \sqcup T) \to C_i(S) \otimes C_i(T)$. The corresponding natural isomorphism for $C_1 \otimes C_2$ is 
    \[
    \begin{tikzcd}
        (C_1 \otimes C_2)(S \sqcup T) =
         C_1(S \sqcup T) \otimes C_2(S \sqcup T) \arrow[d,"\eta^1_{S,T} \otimes \eta^2_{S,T}"]  \\
         C_1(S) \otimes C_1(T) \otimes C_2(S) \otimes C_2(T) \arrow[d,"\id{} \otimes \text{flip} \otimes \id{}"] \\
         (C_1 \otimes C_2)(S) \otimes (C_1 \otimes C_2)(T) = C_1(S) \otimes C_2(S) \otimes C_1(T) \otimes C_2(T)
    \end{tikzcd}    
    \]
    where flip denotes the symmetry of $\GCStar$ (taking the grading into account).
\end{definition}

Our main example of a coefficient system is constructed as follows: let $V$ be a finite-dimensional complex inner product space and let $\rho \colon G \to U(V)$ be a unitary representation of $G$. Let
\begin{equation} \label{eqn:D}
    D = \Endo{V}^{\otimes \infty}
\end{equation}
be the infinite UHF-algebra associated to $V$. The infinite tensor product of the adjoint action of $G$ on $\Endo{V}$ turns this into a $G$-$C^*$-algebra, which we consider to be trivially graded. 

To define the stabilisation of $D$, we need the compact operators, which we turn into a $G$-$C^*$-algebra as follows: let $H_0$ be an infinite-dimensional separable Hilbert space and define $H_G = \ell^2(G) \otimes H_0$, where $\ell^2(G)$ is the finite-dimensional Hilbert space given by the direct sum over all irreducible representations of $G$. Let $\bK = \bK(H_G)$ be the compact operators on $H_G$. This is a $G$-$C^*$-algebra where $G$ acts on $\bK(\ell^2(G) \otimes H_0)$ via the adjoint action. We consider $D \otimes \bK$ equipped with the tensor product action of $G$, which we denote by $\alpha$, and the trivial grading.

\begin{definition}
    Let $C$ be a coefficient system. We call it \emph{stable} if for every finite set $S$ we have $C(S) \otimes \bK \cong C(S)$ as graded $G$-$C^*$-algebras. 
\end{definition}

To define the ring spectrum we will need two non-trivially graded $C^*$-algebras that feature in the definition of $KU^D$. The first one is the graded suspension algebra $\grS = C_0(\R)$ equipped with the grading by odd and even functions. This is not only an algebra, but also a coassociative and cocommutative coalgebra with respect to the comultiplication
\(
    \Delta \colon \grS \to \grS \otimes \grS \ ,
\)
where the tensor product in the codomain is the graded one. 

The second $\Z/2\Z$-graded $C^*$-algebra needed in the construction is the complex Clifford algebra. Let $S \in \Sigma$ be a finite set and let $\C[S]$ be the finite-dimensional complex inner product space equipped with the symmetric bilinear form 
\[
    b(v, w) = \sum_{s \in S} v_s w_s\ ,
\]
where $v_s, w_s$ denote the components with respect to the orthonormal basis~$S$. The Clifford algebra $\Cliff{S}$ is defined to be the quotient
\[
    \Cliff{S} = T(\C[S])/I(\C[S],b)
\]
where $T(W) = \bigoplus_{n \in \N_0} W^{\otimes n}$ is the tensor algebra and $I(W,b)$ is the ideal generated by $v \otimes w - w \otimes v - 2b(v,w)\,1$ for $v,w \in W$ where $b$ is a bilinear form on $W$. To define the grading we declare the elements $w \in W$ to be odd. There is a natural isomorphism of graded algebras
\[
    \Cliff{S \sqcup T} \cong \Cliff{S} \otimes \Cliff{T}\ ,
\]
which turns $S \mapsto \Cliff{S}$ into a coefficient system with trivial $G$-action. For $n \in \N$ let $\Cln{n} = \Cliff{\{1,\dots,n\}}$. We also define 
\[
    (D \otimes \bK)(S) = \bigotimes_{s \in S} D \otimes \bK
\]
and note the natural isomorphism
\[
    (D \otimes \bK)(S \sqcup T) = \bigotimes_{s \in S \sqcup T} D \otimes \bK \cong \left(\bigotimes_{s \in S} D \otimes \bK\right) \otimes \left(\bigotimes_{s \in T} D \otimes \bK\right)
\]
given by reordering the tensor factors. Equipped with this natural isomorphism $S \mapsto (D \otimes \bK)(S)$ is a trivially graded stable coefficient system, where the group acts on $D$ and $\bK$ as described above.

In order to describe Bott periodicity in the spectral picture we define $\hat{\eta}_1 \in \hom_{\text{gr}}(\grS, C_0(\R) \otimes \Cln{1})$ to be the Bott element as in \cite[Def.~1.26]{book:HigsonGuentner}. Taking the adjoint map we can view this element as a base-point preserving continuous map 
\[
    \tilde{\eta}_1 \colon S^1 \to \hom_{\text{gr}}(\grS, \Cln{1})
\]
and extend it to 
\[
    \eta_1 \colon S^1 \to \hom_{\text{gr}}(\grS, \Cln{1} \otimes D \otimes \bK) \quad , \quad z \mapsto \tilde{\eta}_1(z) \otimes 1_D \otimes e\ .
\]
Note that $1_D \otimes e$ is fixed by the group action on $D \otimes \bK$. Now we define the spectrum $KU^D$ as follows:

\begin{definition} \label{def:KUD}
Let $G$ be a finite group, $V$ be a unitary $G$-re\-pre\-sen\-tation and let $D$ be the associated UHF $G$-$C^*$-algebra as defined in \eqref{eqn:D}. Let $(KU^D_n)_{n \in \N_0}$ denote the sequence of $G$-spaces 
\[ 
    KU^D_{n}=\hom_{\text{gr}}(\grS,\, (\C\ell \otimes D \otimes \bK)(\{1,\dots,n\})\,) \ ,
\]
where $S \mapsto (\C\ell \otimes D \otimes \bK)(S)$ is the coefficient system defined above. Notice that there is a natural $(\Sigma_n \times G)$-action on $(\C\ell \otimes D \otimes \bK)(\{1,\dots,n\})$ by functoriality. The spaces $KU^D_n$ can be equipped with a multiplication as follows
\[
    \mu_{m,n} \colon KU^D_m \wedge KU^D_n \to KU^D_{m+n} \quad, \quad \varphi \wedge \psi \mapsto (\varphi \otimes \psi) \circ \Delta\ .
\]
There are corresponding unit maps $\eta_n \colon S^n \to KU^D_n$ defined inductively by 
\(
    \eta_n = (\eta_{n-1} \otimes \eta_1) \circ \Delta
\)
with $\eta_1$ as above.
\end{definition}

\begin{remark}
    Non-equivariant versions of the spectrum $KU^{\cD}$ have been defined in \cite{paper:DP-Units} for all strongly self-absorbing $C^*$-algebras $\cD$. In this case the spectrum $KU^{\cD}$ represents topological $K$-theory with coefficients in $K_0(\cD)$, i.e. a localisation of $K$-theory. For instance, 
    \[KU^{\C} \simeq KU^{\mathcal{Z}} \simeq KU^{\mathcal{O}_{\infty}} \simeq KU \qquad \text{ and } \qquad KU^{M_n^{\infty}} \simeq KU\!\left[\tfrac{1}{n}\right] \]
    where $KU$ denotes the usual complex topological $K$-theory spectrum.
\end{remark}

 Given two graded $G$-$C^*$-algebras $A$ and $B$ denote by $[A,B]_G$ the homotopy classes of grading preserving $*$-homomorphisms $A \to B$. Let $\hat{\C}^2$ be the graded Hilbert space with $\hat{\C}^2_+ = \C = \hat{\C}^2_-$ and let $\hat{H}_G = \hat{\C}^2 \otimes H_G = \hat{\C}^2 \otimes L^2(G) \otimes H_0$. Let $\widehat{\bK} = \bK(\hat{H}_G)$. This is also a graded $G$-$C^*$-algebra. The following lemma is an equivariant generalisation of the converse functional calculus in \cite{paper:Trout}. In fact, the equivariant case can be reduced to the non-equivariant one via the Green-Julg theorem. A proof along these lines (which works for the much more general setting of proper locally compact groupoids with Haar system) can be found for example in \cite[Prop.~3.1]{paper:FuncEquivKSpec}. 
\begin{lemma} \label{lem:htpy_sets}
    Let $G$ be a finite group and let $B$ be a graded $G$-$C^*$-algebra. Then the converse functional calculus explained in \cite[Thm.~4.7]{paper:Trout} provides an isomorphism of groups 
    \[
        [\grS, B \otimes \widehat{\bK}]^G \cong KK_G(\C,B)\ ,
    \]
    which is natural in $B$. 
\end{lemma}

We are now in the position to prove the main theorem of this section.

\begin{theorem} \label{KU}
The spaces $KU^D_{*}$ together with the maps $\mu_{*,*}$ and $\eta_{*}$ form a commutative symmetric $G$-ring spectrum (in the sense of \cite[Def.~2.2]{paper:Hausmann} for the trivial universe), with coefficients (for $H \subseteq G)$
\[ 
    \pi_{n}^{H}(KU^D) \cong K_{n}^{H}(D)\ .
\]
\end{theorem}

\begin{proof}
The proof that the sequence of spaces $(KU^D_n)_{n \in \N_0}$ forms a commutative symmetric $G$-ring spectrum is very similar to the non-equivariant case discussed in \cite[Thm.~4.2]{paper:DP-Units}. Therefore we omit the details and just point out that $KU^D_n$ comes equipped with a $G$-action, which is trivial everywhere except on $D \otimes \bK$. In particular, it commutes with the $\Sigma_n$-action.

To compute $\pi_n^H(KU^D)$ we need to determine
\(
    [S^{n + k}, KU^D_k]^H
\) for $k \in \N_0$. 
Observe that 
\[
    KU^D_k \cong \hom_{\text{gr}}\left(\grS, \Cln{k} \otimes (D \otimes \bK)^{\otimes k}\right) 
\]
and therefore 
\begin{equation} \label{eqn:htpy_groups}
    \left[S^{n +k}, KU^D_k\right]^H \cong \left[\grS, C_0(\R^n) \otimes C_0(\R^k, \Cln{k}) \otimes (D\otimes\bK)^{\otimes k}\right]^H\ .
\end{equation}
From \cite[Cor.~14.5.3]{book:Blackadar} we deduce the following isomorphisms of graded $G$-$C^*$-algebras:
\[
    \Cln{k} \otimes \widehat{\bK} \cong \Cln{k-1} \otimes \Cln{1} \otimes \widehat{\bK} \cong \Cln{k-1} \otimes \Cln{1} \otimes \bK \cong \Cln{k} \otimes \bK\ ,
\]
where $\bK$ on the right hand side denotes the trivially graded compact operators. By Lemma~\ref{lem:htpy_sets} we therefore can express the right hand side of \eqref{eqn:htpy_groups} as equivariant $K$-theory groups as follows:
\begin{align*}
    & KK_H(\C, C_0(\R^n) \otimes C_0(\R^k, \Cln{k}) \otimes (D\otimes\bK)^{\otimes k} ) \\
    \cong\ & KK_H(\C, C_0(\R^n) \otimes (D\otimes\bK)^{\otimes k} ) \cong K_n^H( (D\otimes \bK)^{\otimes k} )\ ,
\end{align*}
where the first isomorphism is a consequence of Bott periodicity in equivariant $KK$-theory. For $m \geq k$ the above identification fits into the following commutative diagram
\[
    \begin{tikzcd}
        \left[S^{n + k}, KU^D_k\right]^H \arrow[d] \arrow[r,"\cong"] & K_n^H( (D\otimes \bK)^{\otimes k} ) \arrow[d] \\
        \left[S^{n + m}, KU^D_m\right]^H \arrow[r,"\cong" below] & K_n^H( (D\otimes \bK)^{\otimes m} )         
    \end{tikzcd}
\]
where the right hand vertical arrow is induced by inclusions $(D \otimes \bK)^{\otimes k} \to (D \otimes \bK)^{\otimes m}$ that introduce tensor factors $1 \otimes e$. All of these maps are isomorphisms in $K$-theory. Hence, 
\[
    \pi_n^H(KU^D) = \colim_{k \to \infty} \left[S^{n + k}, KU^D_k\right]^H \cong K_n^H(D)\ . \qedhere
\]
\end{proof}

\begin{remark}
With only minor modifications one can adopt the above proof to see that $KU^D_n$ extends to a commutative symmetric $G$-ring spectrum for an arbitrary choice of $G$-universe in the sense of \cite[Def.~2.2]{paper:Hausmann}. We omit the details here, since we will only make use of $KU^D$ as a naive $G$-spectrum. 
\end{remark}

\subsection{Equivariant units of $KU^D$}
\label{equnits}

Given a (commutative) symmetric $G$-ring spectrum $R$ with unit maps $\eta_n$ and multiplication maps $\mu_{m,n}$, %the assignment
\[\textbf{n} \mapsto (G/H \mapsto (\Omega^{n}R_{n})^H)\]
defines a (commutative) $\cI$-$O_G$-monoid $\Omega^{\bullet} R$. 
A morphism $\alpha: \textbf{m} \rightarrow \textbf{n}$ in~$\cI$ induces a unique permutation $\bar{\alpha}: \textbf{n} = \textbf{l} \sqcup \textbf{m} \rightarrow \textbf{n}$, which acts like $\alpha$ on the last $m$ entries and is monotonic on the first $l$ entries. This allows to define a map $(\Omega^{m}R_{m})^H \rightarrow (\Omega^{n}R_{n})^H$ by sending $f \in (\Omega^{m}R_{m})^H$ to the composition
\[S^{n} \xrightarrow{\bar{\alpha}^{-1}} S^{n}=S^{l} \wedge S^{m} \xrightarrow{\eta_{l} \wedge f} R_{l} \wedge R_{m} \xrightarrow{\mu_{l,m}} R_{n} \xrightarrow{\bar{\alpha}} R_{n}\ ,\]
thus yielding functoriality of $\Omega^{\bullet} R$ with respect to injective maps.

Non-equivariantly, the group of units of $R$ is classically defined by the following pullback (see \cite[Def.~11.3]{paper:Lind}):
% https://q.uiver.app/#q=WzAsNCxbMCwwLCJcXHRleHR7R0x9XzFSIl0sWzAsMSwiXFx0ZXh0e0dMfV8xKFxccGlfMChSKSkiXSxbMSwwLCJcXE9tZWdhXntcXGluZnR5fVIiXSxbMSwxLCJcXHBpXzAoUikiXSxbMiwzXSxbMywxLCIiLDIseyJzdHlsZSI6eyJib2R5Ijp7Im5hbWUiOiJub25lIn0sImhlYWQiOnsibmFtZSI6Im5vbmUifX19XSxbMCwxXSxbMCwyLCIiLDEseyJzdHlsZSI6eyJ0YWlsIjp7Im5hbWUiOiJob29rIiwic2lkZSI6InRvcCJ9fX1dLFsxLDMsIiIsMSx7InN0eWxlIjp7InRhaWwiOnsibmFtZSI6Imhvb2siLCJzaWRlIjoidG9wIn19fV1d
\begin{equation} 
\label{eqn:GL1_diag}
\begin{tikzcd}
	{\text{GL}_1R} & {\Omega^{\infty}R} \\
	{\text{GL}_1(\pi_0(R))} & {\pi_0(R)}
	\arrow[hook, from=1-1, to=1-2]
	\arrow[from=1-1, to=2-1]
	\arrow[from=1-2, to=2-2]
	\arrow[hook, from=2-1, to=2-2]
\end{tikzcd}
\end{equation}

In the equivariant setting we want this condition to be realised for all fixed points at the same time.
Hence we define the $\cI$-$O_G$-\textit{space of units} of~$R$ (which we denote by GL$_1^{\bullet}R$ again) by the same formal pullback \eqref{eqn:GL1_diag} as above, but in the category of $\cI$-$O_G$-spaces. 
Unraveling the definition shows that it encodes the information about invertible elements for all fixed points, namely $(\text{GL}_1^{\bullet}R(\textbf{n}))(G/H)$ is made up of those elements of $(\Omega^nR_n)^H$ which are sent to units on $\pi_0$, where $\pi_0$ is a $O_G$-set, defined by $\pi_0(X)(G/H)=\pi_0^H(X)$. For semistable commutative symmetric $G$-ring spectra this has the correct homotopy type (see \cite[Rem.~11.4]{paper:Lind}), so in particular for $KU^D$. The machinery illustrated in Section \ref{prel_1} allows us to define the special $\Gamma$-$G$-space of units $H_{\cI}\text{GL}_1^{\bullet}R$ and its associated spectrum.

\begin{definition} \label{def:spec_units}
The $G$-\textit{spectrum of units} of a commutative symmetric $G$-ring spectrum~$R$ is
\[gl_{1}R=EH_{\cI}\text{GL}_1^{\bullet}R\ .\]
\end{definition}

\begin{remark}
    A similar approach to units of equivariant ring spectra has been worked out in \cite[Sec.~9]{paper:Santhanam}, with the difference that the author uses special $\Gamma$-$G$-spaces (actually a slight variation, namely $\Gamma_G$-spaces) directly.
    Let us show how for a special $\Gamma$-$O_G$-space arising from a stable commutative $\cI$-$O_G$-monoid $\Omega^{\bullet}R$ for a symmetric commutative $G$-ring spectrum $R$ the two machines define weakly equivalent objects. We will denote the $\Gamma$-$O_G$-space of units defined in \cite{paper:Santhanam} by $UR$. It is defined by the  pullback 
   % https://q.uiver.app/#q=WzAsNCxbMCwwLCJcXHdpZGV0aWxkZXtcXHRleHR7R0x9XzF9XntcXGJ1bGxldH1SIl0sWzAsMSwiXFx0ZXh0e0dMfV8xKFxccGlfMChIX3tcXG1hdGhjYWx7SX19XFxPbWVnYV57XFxidWxsZXR9UikpIl0sWzEsMCwiSF97XFxtYXRoY2Fse0l9fVxcT21lZ2Fee1xcYnVsbGV0fVIiXSxbMSwxLCJcXHBpXzAoSF97XFxtYXRoY2Fse0l9fVxcT21lZ2Fee1xcYnVsbGV0fVIpIl0sWzIsM10sWzMsMSwiIiwyLHsic3R5bGUiOnsiYm9keSI6eyJuYW1lIjoibm9uZSJ9LCJoZWFkIjp7Im5hbWUiOiJub25lIn19fV0sWzAsMV0sWzAsMiwiIiwxLHsic3R5bGUiOnsidGFpbCI6eyJuYW1lIjoiaG9vayIsInNpZGUiOiJ0b3AifX19XSxbMSwzLCIiLDEseyJzdHlsZSI6eyJ0YWlsIjp7Im5hbWUiOiJob29rIiwic2lkZSI6InRvcCJ9fX1dXQ==
\[\begin{tikzcd}
	{UR} & {H_{\mathcal{I}}\Omega^{\bullet}R} \\
	{\text{GL}_1(\pi_0(H_{\mathcal{I}}\Omega^{\bullet}R))} & {\pi_0(H_{\mathcal{I}}\Omega^{\bullet}R)}
	\arrow[hook, from=1-1, to=1-2]
	\arrow[from=1-1, to=2-1]
	\arrow[from=1-2, to=2-2]
	\arrow[hook, from=2-1, to=2-2]
\end{tikzcd}\]
where $H_{\mathcal{I}}\Omega^{\bullet}R$ is the special $\Gamma$-$O_G$-space associated to $\Omega^{\bullet}R$. Note that $H_{\cI}GL_1^\bullet R$ also fits into a homotopy pullback diagram
\[
    \begin{tikzcd}
        H_{\cI}GL_1^\bullet R \ar[r] \ar[d] & H_{\cI} \ar[d] \Omega^\bullet R \\
        \pi_0(H_{\cI}GL_1^\bullet R) \ar[r] & \pi_0(H_{\cI}\Omega^\bullet R)
    \end{tikzcd}
\]
The map of $\Gamma$-$G$-spaces $H_\cI GL_1^\bullet R \to H_\cI \Omega^{\bullet} R$ is an isomorphism on all higher homotopy groups. Hence, it suffices to show that
\[
\pi_0(UR(\textbf{1}^+)) \cong \pi_0((H_{\cI}GL_1^\bullet R)(\mathbf{1}^+)) \ .
\] 
By stability we have $(H_{\cI}\text{GL}_1^{\bullet}R)(\textbf{1}^+) = \text{GL}_1^{\bullet}R_{h\cI} \simeq \text{GL}_1^{\bullet}R(\textbf{1})$ and therefore
\[\pi_0((H_{\cI}\text{GL}_1^{\bullet}R)(\textbf{1}^+)) \cong 
\pi_0(\text{GL}_1^{\bullet}R(\textbf{1})) \cong \text{GL}_1(\pi_0(\Omega^{\infty} R))\ .
\]
On the other hand, since $(H_{\cI}\Omega^{\bullet}R)(\textbf{1}^+)=\Omega^{\bullet}R_{h\cI} \simeq \Omega^{\infty}R$ we have that $UR(\textbf{1}^+)$ is defined by
% https://q.uiver.app/#q=WzAsNCxbMCwwLCJcXHdpZGV0aWxkZXtcXHRleHR7R0x9XzF9XntcXGJ1bGxldH1SKFxcdGV4dGJmezF9KSJdLFswLDEsIlxcdGV4dHtHTH1fMShcXHBpXzAoXFxPbWVnYV57XFxpbmZ0eX1SKSkiXSxbMSwwLCJcXE9tZWdhXntcXGluZnR5fVIiXSxbMSwxLCJcXHBpXzAoXFxPbWVnYV57XFxpbmZ0eX1SKSJdLFsyLDNdLFszLDEsIiIsMix7InN0eWxlIjp7ImJvZHkiOnsibmFtZSI6Im5vbmUifSwiaGVhZCI6eyJuYW1lIjoibm9uZSJ9fX1dLFswLDFdLFswLDIsIiIsMSx7InN0eWxlIjp7InRhaWwiOnsibmFtZSI6Imhvb2siLCJzaWRlIjoidG9wIn19fV0sWzEsMywiIiwxLHsic3R5bGUiOnsidGFpbCI6eyJuYW1lIjoiaG9vayIsInNpZGUiOiJ0b3AifX19XV0=
\[\begin{tikzcd}
	{UR(\textbf{1}^+)} & {\Omega^{\bullet}R_{h\cI}} \\
	{\text{GL}_1(\pi_0(\Omega^{\infty}R))} & {\pi_0(\Omega^{\infty}R)}
	\arrow[hook, from=1-1, to=1-2]
	\arrow[from=1-1, to=2-1]
	\arrow[from=1-2, to=2-2]
	\arrow[hook, from=2-1, to=2-2]
\end{tikzcd}\]
Hence by definition 
\[
\pi_0(UR(\textbf{1}^+)) \cong \text{GL}_1(\pi_0(\Omega^{\infty}R)) \cong \pi_0(H_{\cI}\text{GL}_1^{\bullet}R(\textbf{1}^+))\ .
\]
\end{remark} 

We will focus on the equivariant units of the $G$-ring spectrum $R=KU^D$ constructed in Section \ref{spectrum}. The following result is basically a consequence of Bott periodicity.

\begin{lemma}
The $\cI$-$O_G$-space $\text{GL}_1^{\bullet}KU^D$ is stable.
\end{lemma}

\begin{proof}
We need to check that all morphisms $\mathbf{m} \to \mathbf{n}$ in~$\cI$ with $m > 0$ are mapped to weak equivalences for all cosets $G/H$. Let $H \subseteq G$. Note that for $\pi_0$ we have
\[ \pi_{0} ((\text{GL}_1^{\bullet}KU^D)(\textbf{m}))(G/H) \cong \pi_{0}((\text{GL}_1^{\bullet}KU^D)(\textbf{n}))(G/H) \cong \text{GL}_1(K_0^H(D))\]
for all \textbf{m}, \textbf{n}, whereas for $k > 0$ we have the following commutative diagram
% https://q.uiver.app/#q=WzAsNCxbMCwwLCJcXHBpX2soXFx0ZXh0e0dMfV8xXntcXGJ1bGxldH1SKShcXHRleHRiZnttfSkoRy9IKSJdLFswLDEsIlxccGlfa1xcT21lZ2Fee219S1VeRF9tIl0sWzEsMCwiXFxwaV9rKFxcdGV4dHtHTH1fMV57XFxidWxsZXR9UikoXFx0ZXh0YmZ7bn0pKEcvSCkiXSxbMSwxLCJcXHBpX2tcXE9tZWdhXntufUtVXkRfbiJdLFsyLDMsIlxcY29uZyJdLFszLDEsIiIsMix7InN0eWxlIjp7ImJvZHkiOnsibmFtZSI6Im5vbmUifSwiaGVhZCI6eyJuYW1lIjoibm9uZSJ9fX1dLFswLDEsIlxcY29uZyIsMl0sWzAsMl0sWzEsMywiXFxjb25nIiwyXV0=
\[\begin{tikzcd}
	{\pi_k((\text{GL}_1^{\bullet}KU^D))(\textbf{m})(G/H)} & {\pi_k((\text{GL}_1^{\bullet}KU^D)(\textbf{n}))(G/H)} \\
	{\pi_k(\Omega^{m}KU^D_m)} & {\pi_k(\Omega^{n}KU^D_n)}
	\arrow["\cong", from=1-2, to=2-2]
	\arrow["\cong"', from=1-1, to=2-1]
	\arrow[from=1-1, to=1-2]
	\arrow[from=2-1, to=2-2]
\end{tikzcd}\]
where the bottom horizontal arrow is an isomorphism when $m > 0$.
\end{proof}


\section{Automorphisms of $D \otimes \bK$}

\subsection{The $\cI$-$G$-space $\cG_D$}
\label{G_D}

The $C^*$-algebra $D=\Endo{V}^{\infty}$ is unital, separable and strongly self-absorbing. It has been shown in \cite{paper:DP-Dixmier} that the automorphism group of the stabilisation of such algebras carries an infinite loop space structure.
In order to extend this to the equivariant case we consider $D \otimes \bK$ as a $G$-$C^*$-algebra for a finite group $G$ (which we will later choose to be $\Z/p\Z$) with the action on $D$ defined in \eqref{eqn:D}. 

Let $\cG_D(\textbf{n})=\Aut{(D \otimes \bK)^{\otimes n}}$ (with $(D \otimes \bK)^{\otimes 0} :=\C$, so that $\cG_D(\textbf{0})$ is the trivial group).
The functor $\cG_D$ is defined on a morphism $\alpha \colon \textbf{m} \rightarrow \textbf{n}$ in~$\cI$ as $\cG_D(\alpha)(g)=\bar{\alpha} \circ (\id{(D \otimes \bK)^{\otimes (n-m)}} \otimes g) \circ \bar{\alpha}^{-1}$, with $g \in \cG_D(\textbf{m})$ and $\bar{\alpha}$ as explained at the beginning of Section \ref{equnits}. The tensor product equips $\cG_D(\textbf{n})$ with the structure of a commutative $\cI$-$G$-monoid, via
\[ \mu_{m,n}: \cG_D(\textbf{m}) \times \cG_D(\textbf{n}) \rightarrow \cG_D(\textbf{m} \sqcup \textbf{n})\ , \qquad (g,h) \mapsto g \otimes h\ .\] 
Note that in this case $\cG_D(\textbf{n})$ has both a group multiplication and a monoid structure and they satisfy the Eckmann-Hilton condition from \cite[Def.~3.1]{paper:DP-Units}. It will prove useful to have a compatibility condition between the two.

\begin{definition}
    Let $X$ be a $\cI$-$G$-monoid such that $X(\textbf{n})$ is also a group for each \textbf{n}. Let $\iota_m:\textbf{0} \rightarrow \textbf{m}$ be the unique such morphism in $\cI$. We say that $X$ has \textit{compatible inverses} if there is a path from $(\iota_m \sqcup \id{\textbf{m}})_*(g) \in X(\textbf{m} \sqcup \textbf{m})$ to $(\id{\textbf{m}} \sqcup \iota_m)_*(g)$ for all $m$ and $g \in X(\textbf{m})$.
\end{definition}

If $X$ has compatible inverses, then for each $g \in X(\textbf{m})$ there is a path connecting $\mu_{m,m}(g,g^{-1}) \in X(\textbf{m} \sqcup \textbf{m})$ and $1_{m \sqcup m} \in X(\textbf{m} \sqcup \textbf{m})$ \cite[Lem. 3.3]{paper:DP-Units}. 

From now on we restrict to $G = \Z/p\Z$ for a prime $p \in \N$.
\begin{lemma}
\label{GD_stable}
    $\cG_D$ is a stable commutative $\cI$-$\Z/p\Z$-monoid with compatible inverses.
\end{lemma}

\begin{proof}
    See \cite[Thm. 4.5]{paper:DP-Units} for $H=\{e\}$. The proof for $H = \Z/p\Z$ is an easy adaptation of \cite[Lem. 4.2]{paper:EP-Circle} where this is proven for $H = \bT$. 
\end{proof}

Recall we can always regard an $\cI$-$G$-monoid as an $\cI$-$O_G$-monoid by taking fixed points, since this preserves monoidality. By a slight abuse of notation we will refer to the latter again as $\cG_D$ in Section \ref{final_2}, where we will compare it with the $\cI$-$O_G$-space of units.


\subsection{$\Z/p\Z$-equivariant homotopy type of $\Aut{D \otimes \bK}$}
\label{Aut}

Let $G = \Z/p\Z$. Since the $G$-$C^*$-algebra $D$ defined in \eqref{eqn:D} is an infinite tensor product of a single representation, one can compute its equivariant $K$-theory using continuity. Note that the map $T \mapsto T \otimes 1_D$ induces multiplication by $V$ on $K$-theory. Therefore,
\[
    K_0^{G}(\C) \cong K_0^G(\Endo{V}^{\otimes n}) \cong \Z[t]/(t^p - 1) =: R \quad \text{and} \quad K_0^{G}(D) \cong R[p_V^{-1}]\ ,
\]
where $p_V \in R$ is the polynomial corresponding to the $G$-representation $V$. The group $K_0^G(D)$ is in fact a ring. But it also has the structure of an ordered group. The positive cone of $K_0^G(D)$ corresponds under this isomorphism to $R[p_V^{-1}]_+$, which is defined for $q \in R$ and $k \in \N_0$ by
\[
    \frac{q}{p_V^k} \in R[p_V^{-1}]_+ \  \Leftrightarrow \  q\cdot p_V^l \text{ has non-negative coefficients for some } l \in \N_0\ .
\]
Taking this order structure into account we need to consider the subgroup of positive units: Let $R_+ \subset R$ be the semiring of positive elements contained in a commutative unital ring $R$ with an ordered group structure and define
%For a semiring $R_+ \subset R$ contained in a commutative unital ring $R$, define
\[
    GL_1(R)_+= \{ x \in GL_1(R) \cap R_+ \ | \ x^{-1} \in R_+ \}\ ,
\]
i.e.\ those elements of $R_+$ which are invertible as elements of $R$ and whose inverse is in $R_+$ (see Remark~\ref{rem:glitch} for the reason why this definition is different from the one in \cite[p.~17]{paper:EP-Circle}). 

As we have already seen above, some of the results from \cite{paper:EP-Circle} are easily adapted to the subgroup $\Z/p\Z \subseteq \bT$. For example, determining the equivariant homotopy type of $\Aut{D \otimes \bK}$ can be achieved using the same approach as in \cite{paper:EP-Circle}. 

\begin{lemma} \label{lem:contr_htpies}
Let $e \in \bK$ be a projection of rank one that is fixed by the $\Z/p\Z$-action. There are two continuous paths 
\begin{align*}
    \gamma \colon [0,1] &\to \hom_{\Z/p\Z}(D, D \otimes D) \ , \\
    \beta \colon [0,1] &\to \hom_{\Z/p\Z}(D \otimes \bK, (D \otimes \bK)^{\otimes 2})
\end{align*}
(where the equivariant homomorphisms are equipped with the point-norm topology) with the following properties:
\begin{enumerate}[i)]
    \item $\gamma(0)(d) = d \otimes 1_D$, $\gamma(1)(d) = 1_D \otimes d$,
        $\beta(0)(a) = a \otimes (1_D \otimes e)$ and $\beta(1)(a) = (1_D \otimes e) \otimes a$,
    \item $\gamma$ restricts to $(0,1) \to \mathrm{Iso}_{\Z/p\Z}(D, D \otimes D)$ and $\beta$ restricts to $(0,1) \to \mathrm{Iso}_{\Z/p\Z}(D \otimes \bK, (D \otimes \bK)^{\otimes 2})$,
    \item $\beta(t)(1_D \otimes e) = (1_D \otimes e)^{\otimes 2}$ for all $t \in [0,1]$.
\end{enumerate}
\end{lemma}

\begin{proof}
    The path $\gamma$ is the result of patching together two paths that are constructed exactly as in \cite[Lem.~2.3]{paper:EP-Circle} (see also \cite[p.~11]{paper:EP-Circle}). Since the analysis in \cite[Thm.~2.4]{paper:EP-Circle} carries over to $\Z/p\Z$, the stabiliser group $\AutSt{\Z/p\Z,e}{\bK}$ is contractible in the point-norm topology (so in particular path-connected). Therefore the argument in \cite[Lem.~2.5]{paper:EP-Circle} works for $\Z/p\Z$ as well giving a continuous path
    \[
        \overline{\beta} \colon [0,1] \to \hom_{\Z/p\Z}(\bK, \bK \otimes \bK)\ .
    \]
    with properties analogous to the ones listed in the lemma. The path $\beta$ is then obtained by combining $\gamma$ and $\overline{\beta}$.
\end{proof}

\begin{corollary} \label{cor:Aut_contractible}
    Let $e \in \bK$ be a projection of rank one fixed by $\Z/p\Z$. The group $\eqAut{\Z/p\Z}{D}$ and the stabiliser subgroup $\AutSt{\Z/p\Z, 1_D\otimes e}{D \otimes \bK}$ fixing the projection $1_D \otimes e$ are both contractible.
\end{corollary}
\begin{proof}
    This is proven exactly as in \cite[Thm.~2.6]{paper:EP-Circle}. The contracting homotopies are constructed from the paths $\gamma$ and $\beta$, respectively, given in Lemma~\ref{lem:contr_htpies}.
\end{proof}

The above corollary allows us to determine the equivariant homotopy type of $\Aut{D \otimes \bK}$ using the evaluation map
\begin{equation} \label{eqn:Aut_to_proj}
    \eqAut{H}{D \otimes \bK} \to \Proj{}{(D \otimes \bK)^H} \quad, \quad \beta \mapsto \beta(1_D \otimes e)
\end{equation}
suitably restricted in the codomain for $H \subseteq \Z/p\Z$. As in \cite[Lem.~2.8]{paper:EP-Circle} Rieffel's criterion applies to the $\Z/p\Z$-action on $\bK$ and since colimits of saturated actions are saturated, it follows that the inclusion map induces an isomorphism
\[
    K_0((D \otimes \bK)^{\Z/p\Z}) \to K_0^{\Z/p\Z}(D \otimes \bK) \cong K_0^{\Z/p\Z}(D) \ .
\]
By stability we also have $K_0(D \otimes \bK) \cong K_0(D)$. Combining \eqref{eqn:Aut_to_proj} with the natural map $\pi_0(\Proj{}{(D \otimes \bK)^H}) \to K_0((D \otimes \bK)^H)$ and the above isomorphisms therefore gives a map
\[
   \kappa \colon \pi_0(\eqAut{H}{D \otimes \bK}) \to K_0^{H}(D)\ .
\]

\begin{lemma}
\label{units+}
The map $\kappa$ restricts to a group isomorphism
\[ \pi_{0}(\eqAut{H}{D \otimes \bK}) \xrightarrow{\cong} GL_1(K_{0}^H(D))_+\ .\]
\end{lemma}

\begin{proof}
The case $H = \{e\}$ is proven in \cite[Cor. 2.17]{paper:DP-Dixmier}. That $\kappa$ is multiplicative in the case $H = \Z/p\Z$ is shown as in the proof of \cite[Lem.~2.9]{paper:EP-Circle} and is based on the paths constructed in Lemma~\ref{lem:contr_htpies}. The argument also proves that $[\alpha^{-1}] \in \pi_0(\eqAut{\Z/p\Z}{D \otimes \bK})$ is the inverse of $[\alpha]$. Hence, $\kappa([\alpha^{-1}])$ is the inverse of $\kappa([a])$ and in particular positive. Thus, the restriction gives a well-defined group homomorphism to $GL_1(K_0^H(D))_+$. 

Showing its injectivity can be done using the same argument as in the proof of \cite[Lem.~2.9]{paper:EP-Circle}. For reasons that will become clear in Remark~\ref{rem:glitch} we spell out the proof of surjectivity explicitly. Recall that 
\[
    K_0^{\Z/p\Z}(\C) \cong \Z[t]/(t^p - 1) =: R \quad \text{and} \quad K_0^{\Z/p\Z}(D) \cong R[p_V^{-1}]\ ,
\]
where $p_V(t)$ denotes the polynomial corresponding to the defining representation of $D = \Endo{V}^{\otimes \infty}$. Let $q \in GL_1(R[p_V^{-1}])_+$. Then
\[
    q = \frac{r}{p_V^k}
\]
for some $r \in R_+$, $k \in \N_0$ and by our definition of $GL_1(R[p_V^{-1}])_+$ there is $s \in R_+$ such that 
\(
    r \cdot s = p_V^l
\) for some $l \in \N_0$. Since $r$ and $s$ are both positive, this corresponds to a tensor product decomposition of the $\Z/p\Z$-representation $V^{\otimes l}$, i.e.\ we have two representations $V_r$ and $V_s$ corresponding to $r$ and $s$, respectively, such that
\begin{equation} \label{eqn:V_End_iso}
    V_r \otimes V_s \cong V^{\otimes l} \quad \text{and} \quad \Endo{V_r} \otimes \Endo{V_s} \cong \Endo{V}^{\otimes l}\ .
\end{equation}
First assume that $V_s$ contains a copy of the trivial representation. Then $V_r$ is a subrepresentation of $V^{\otimes l}$. Observe that because of \eqref{eqn:V_End_iso} the $G$-$C^*$-algebra~$D$ tensorially absorbs $\Endo{V^r}$. Let  
\[
    \psi_r \colon \Endo{V_r} \otimes D \cong \Endo{V_r} \otimes (\Endo{V_r} \otimes \Endo{V_s})^{\otimes \infty} \to D\ .
\]
be the equivariant $*$-isomorphism that shifts the tensor factors and identifies the result with $D$ using \eqref{eqn:V_End_iso}.
Since $V_r \otimes H \cong H$ for $H = \ell^2(\Z/p\Z) \otimes H_0$, we can also choose an isomorphism 
\[
    \phi_r \colon \bK(H) \otimes \Endo{V_r} \to \bK(H \otimes V_r) \to \bK(H)\ .
\]
The composition $(\phi_r \otimes \id{D}) \circ (\id{\bK} \otimes \psi_r^{-1})$ gives $\alpha_r \in \eqAut{\Z/p\Z}{D \otimes \bK}$ after swapping $D$ and $\bK$. By construction, $\alpha_r(1_D \otimes e)$ is an invariant projection, whose image is the sub\-re\-pre\-sen\-ta\-tion~$V_r$, i.e.\ $[\alpha_r(1_D \otimes e)] = r \in R_+$. Similarly, we can construct $\alpha_{p_V} \in \eqAut{\Z/p\Z}{D \otimes \bK}$ with $[\alpha_{p_V}(1_D \otimes e)] = p_V$. Since $\kappa$ is a group homomorphism, we have $\kappa([(\alpha_{p_V}^{-1})^k \circ \alpha_r]) = q$. 

If $V_s$ does not contain the trivial representation, then there is $m \in \Z/p\Z$ such that $\C_m \otimes V_s$ does, where $\C_m$ is the character corresponding to $m$. The above argument then shows that $t^{-m}q(t)$ is represented by an automorphism $\beta$. As described in the proof of \cite[Lem.~2.9]{paper:EP-Circle} there is an equivariant automorphism $\alpha_t \in \eqAut{\Z/p\Z}{\bK}$ with the property $[\alpha_t(e)] = t \in R_+$ and we have $\kappa([(\id{D} \otimes \alpha_t^m) \circ \beta]) = q$.
\end{proof}

\begin{remark} \label{rem:glitch}
    There is a small mistake in the proof of \cite[Lem.~2.9]{paper:EP-Circle}: The definition of $GL_1(R_+)$, the surjectivity part of the proof of the lemma and the computation of the generators of $GL_1(R_+)$ are actually only correct under the additional assumption that the prime factors of $p_V$ are positive. This is not automatic as the example
    \[
        p_V(t) = 1+t^3 = (1 - t + t^2)(1+t)
    \]
    shows. The above proof of surjectivity also works in the case $G = \bT$ and fixes this gap. Note that there is another reason why the proof of \cite[Lem.~2.9]{paper:EP-Circle} does not carry over verbatim: The ring $\Z[t]/(t^p - 1)$ contains zero divisors and is therefore not a unique factorization domain.
\end{remark}

\begin{prop}
\label{ProjBU}
There is a homotopy equivalence
\[ \Proj{1 \otimes e}{(D \otimes \bK)^H} \simeq BU(D^H)\ .\]
\end{prop}

\begin{proof}
See \cite[Cor. 2.9]{paper:DP-Dixmier} for $H=\{e\}$. The proof of the case $H = \Z/p\Z$ is the same as the one of \cite[Prop. 3.2]{paper:EP-Circle}.
\end{proof}

\begin{remark}
\label{RemProjBU}
Let $V$ be a right Hilbert $A$-module and note Hilbert $A$-module bundles with fibre isomorphic to $V$ over $X$ are up to isomorphism in bijection with $[X, BU(V)]$, where $U(V)$ denotes the unitary group of $V$. This observation allows us to reinterpret Proposition \ref{ProjBU} in terms of a statement about right Hilbert $D^{\Z/p\Z}$-module bundles. In fact, the proof works by constructing a principal $U(D^{\Z/p\Z})$-bundle 
\[EU(D^{\Z/p\Z})=\frac{U(M(D^{\Z/p\Z}))}{U((1-p_0)M(D^{\Z/p\Z})(1-p_0))}\]
over $\Proj{1 \otimes e}{(D \otimes \bK)^{\Z/p\Z}}$ (where we have set $p_0=e \otimes 1 \in H_G \otimes D^{\Z/p\Z}$) and showing that it is contractible. The bundle $EU(D^{\Z/p\Z})$ defines in turn a right Hilbert $D^{\Z/p\Z}$-module bundle
\[\cH_{\text{univ}}=EU(D^{\Z/p\Z})\times_{U(D^{\Z/p\Z})} D^{\Z/p\Z}\]
over $\Proj{p_0}{(D \otimes \bK)^{\Z/p\Z}}$. We can also define naturally another right Hilbert $D^{\Z/p\Z}$-module bundle over the same base space, namely
\[\tilde{\cH}_{\text{univ}}=\{(p,\xi) \in \Proj{p_0}{(D \otimes \bK)^{\Z/p\Z}} \times (H_G \otimes D^{\Z/p\Z}) : p \cdot \xi = \xi\}\ .\]
Then the map 
\[\alpha: \cH_{\text{univ}} \rightarrow \tilde{\cH}_{\text{univ}}\]
defined by $\alpha([u,b])=(u p_0 u^*,u \cdot (e_{0} \otimes b))$, where $e_0 \in H$ is a unit vector in the image of $p_0$, is a morphism of right Hilbert $D^{\Z/p\Z}$-module bundles, with the properties of being an isomorphism on the fibres and making the triangle
% https://q.uiver.app/#q=WzAsMyxbMCwwLCJcXG1hdGhjYWx7SH1fe1xcdGV4dHt1bml2fX0iXSxbMiwwLCJcXHRpbGRle1xcbWF0aGNhbHtIfX1fe1xcdGV4dHt1bml2fX0iXSxbMSwxLCJcXHRleHR7UHJvan1fe3BfMH0oKEQgXFxvdGltZXMgSylee1xcWi9wXFxafSkiXSxbMCwyXSxbMSwyXSxbMCwxLCJcXGFscGhhIl1d
\[\begin{tikzcd}
	{\mathcal{H}_{\text{univ}}} && {\tilde{\mathcal{H}}_{\text{univ}}} \\
	& {\text{Proj}_{p_0}((D \otimes K)^{\Z/p\Z})}
	\arrow[from=1-1, to=2-2]
	\arrow[from=1-3, to=2-2]
	\arrow["\alpha", from=1-1, to=1-3]
\end{tikzcd}\]
commute. Thus it is an isomorphism of right Hilbert $D^{\Z/p\Z}$-module bundles.
\end{remark}

\begin{prop}
\label{AutProj}
There is a homotopy equivalence
\[\eqAutId{H}{D \otimes \bK} \simeq \Proj{1 \otimes e}{(D \otimes \bK)^H}\]
for $H \subseteq \Z/p\Z$.
\end{prop}

\begin{proof}
The proof is based on the observation that the evaluation map
\[
    \eqAutId{H}{D \otimes \bK} \to \Proj{1 \otimes e}{(D \otimes \bK)^H} \quad, \quad \beta \mapsto \beta(1 \otimes e)
\]
is a principal $\AutSt{\Z/p\Z,1 \otimes e}{D \otimes \bK}$-bundle with contractible fibre by Corollary~\ref{cor:Aut_contractible}. For further details see \cite[Cor. 2.9]{paper:DP-Dixmier} for $H=\{e\}$, and \cite[Prop.~3.3]{paper:EP-Circle} where this is discussed for the circle group $\bT$ with a proof that carries over to the case $H=\Z/p\Z$.
\end{proof}

The following elementary consideration will be needed in order to compute the $K_0$-group of the fixed point algebra of $D$.

\begin{lemma}
\label{coefficients}
Let $V$ be a $\Z/p\Z$-representation which contains at least two non-isomorphic irreducible subrepresentations, and let $p_{V}=\sum_{i=0}^{p-1} a_{i} t^{i}$ be the associated character polynomial. Then there exists $N \in \N$ such that the $N$-th power
\[p_{V}^{N}=\sum_{j=0}^{p-1} b_{j} t^{j}\]
has coefficients $b_{j} \ge 2$ for all $j \in \{0,\dots,p-1\}$.
\end{lemma}

\begin{proof}
We know that the character polynomial $p_{V}(t)=\sum_{i=0}^{d} a_{i} t^{i} \in \frac{\Z[t]}{(t^{p}-1)}$ has at least two nonzero coefficients. First, assume one of them is $a_{1}$ and let the other one be $a_{l}$, $l \in \{0,2,\dots,p-1\}$. By taking the square $p_{V}^{2}$ we get $a_2>0$ and the coefficient of $t^{l+1}$ increases by at least $2$, and by iterating the process we get $N \in \N$ satisfying the claim.\\
If $a_{1}=0$, by assumption there exists $k \in \{2,\dots,p-1\}$ such that $a_{k} \ne 0$. Let $r$ be the multiplicative inverse of $k$ mod $p$. Then the first order coefficient of $p_{V}^{r}$ is nonzero and the above argument applies. 
\end{proof}

\begin{remark}
The above lemma is actually the only reason why we work with a prime number $p$. In fact the results in the paper hold for cyclic groups of any order, provided that a suitable representation $V$ (i.e., such that $p_V$ satisfies the claim of Lemma \ref{coefficients}) is chosen.
\end{remark}

\begin{lemma}
\label{kt_fixed}
The $K$-groups of the fixed-point algebra $D^{\Z/p\Z}$ are given by
\[K_0(D^{\Z/p\Z}) \cong R(\Z/p\Z)[p_V^{-1}] \cong K_0^{\Z/p\Z}(D) \text{ and } K_1(D^{\Z/p\Z})=0\ .\]
\end{lemma}

\begin{proof}
The proof works similarly as in \cite[Lem. 3.5]{paper:EP-Circle}. First note that because the group acts on each factor separately, the fixed-point algebra $D^{\Z/p\Z}$ is isomorphic to the direct limit
\[ \End(V)^{\Z/p\Z} \rightarrow \cdots \rightarrow \bigl(\End(V)^{\otimes n}\bigr)^{\Z/p\Z} \rightarrow \bigl(\End(V)^{\otimes (n+1)}\bigr)^{\Z/p\Z} \rightarrow \cdots,\]
where the connecting $*$-homomorphisms are given by $T \mapsto T \otimes 1$. For each factor we have
\[\bigl(\End(V)^{\otimes n}\bigr)^{\Z/p\Z} \cong \hom_{\Z/p\Z}(V^{\otimes n}, V^{\otimes n})\ ,\]
which is in turn isomorphic to a direct sum of matrix algebras. Combining the decomposition of $V$ into character subspaces 
\[V=\bigoplus_{k \in \Z} V_k\ ,\]
where $V_k=\{v \in V \, | \, \rho(z)v=z^k v \, \, \forall \, z \in \Z/p\Z\}$, with Schur's lemma this turns out to be
\[\bigoplus_{k \in \Z}\hom_{\Z/p\Z}(V_k, V_k) \cong \bigoplus_i M_{n_i}(\C) \otimes \hom_{\Z/p\Z}(\C_{m_i},\C_{m_i}) \cong \bigoplus_i M_{n_i} (\C)\]
where $V_{m_i} \cong \C^{n_i} \otimes \C_{m_i}$. This implies that $K_0(\End(V)^{\otimes n})^{\Z/p\Z}$ is the free abelian group generated by the irreducible subrepresentations of $V^{\otimes n}$. If we apply $K_0$ to each term in the direct limit above and identify a representation with its character polynomial, the connecting isomorphism becomes multiplication by $p_V$. Moreover, Lemma \ref{coefficients} implies in particular that $p_V^N$ contains all powers of $t$ with positive coefficients. The claim for $K_0$ follows by combining these two considerations and continuity of $K$-theory. Since $D^{\Z/p\Z}$ is an AF-algebra we also have $K_1(D^{\Z/p\Z})=0$.\\
For equivariant $K$-theory, note that 
\[K_0^{\Z/p\Z}(\End(V)^{\otimes n}) \cong K_0^{\Z/p\Z}(\C) \cong R(\Z/p\Z)\] 
by stability, and then again by continuity $K_0^{\Z/p\Z}(D) \cong R(\Z/p\Z)[p_V^{-1}]$. The description as a direct limit also shows that this is in fact a ring isomorphism.
\end{proof}

We can now compute the homotopy groups of $\eqAut{\Z/p\Z}{D \otimes \bK}$. By Proposition \ref{ProjBU} and Proposition \ref{AutProj}, this boils down to computing the ones of $U(D^{\Z/p\Z})$. Here the fact that the representation ring of $\Z/p\Z$ has finite rank over $\Z$ plays a fundamental role.

\begin{theorem}
\label{homotopyU}
We have $\pi_{2k}(U(D^{\Z/p\Z}))=0$ and
\[\pi_{2k+1}(U(D^{\Z/p\Z})) \cong R(\Z/p\Z)[p_{V}^{-1}] \cong K_{0}^{\Z/p\Z}(D)\]
for $k \in \N_{0}$.
\end{theorem}

\begin{proof}
The proof works similarly as in \cite[Thm. 3.7]{paper:EP-Circle}. Let $n \in \N_{0}$. The group $\pi_{n}(U(D^{\Z/p\Z}))$ is the direct limit of the sequence
\[\cdots \rightarrow \pi_{n}(U(\End(V^{\otimes k})^{\Z/p\Z}) \rightarrow \pi_{n}(U(\End(V^{\otimes k+1}))^{\Z/p\Z}) \rightarrow \cdots\]
where the connecting homomorphism is induced by the map $T \mapsto T \otimes 1$.\\
Consider the representing polynomial for $V$
\[p_{V}= \sum_{i=0}^{p-1} a_{i}t^{i}\ .\] 
The coefficients $a_{i}$ can be interpreted as the multiplicities of the character subspaces $V_{i}$ of $V$, i.e., $a_{i}=\dim(V_{i})$. Hence, if 
\[p_{V}^{k}= \sum_{j=0}^{p-1} b_{j}t^{j}\ ,\] 
then by Schur's lemma (compare Lemma \ref{kt_fixed})
\[\End(V^{\otimes k})^{\Z/p\Z} \cong \bigoplus_{j=0}^{p-1} M_{b_j}(\C) \otimes \End(W_{j})\ ,\]
where $W_{j} \subset V^{\otimes k}$ are the irreducible subrepresentations, which appear with multiplicity $b_{j}$. Hence,
\[U(\End(V^{\otimes k})^{\Z/p\Z}) \cong \prod_{j=0}^{p-1} U(b_{j})\ .\]
By Lemma \ref{coefficients} there exists $N \in \N$ such that $p_{V}^{N}$ has all coefficients $\ge$ 2. Since the sequence
\[\End(V^{\otimes 0})^{\Z/p\Z} \rightarrow \End(V^{\otimes N})^{\Z/p\Z} \rightarrow \cdots \rightarrow \End(V^{\otimes mN})^{\Z/p\Z} \rightarrow \cdots\]
also has direct limit $D^{\Z/p\Z}$, we can work with $V^{\otimes N}$ instead of $V$ and assume without loss of generality that $p_{V}$ has all coefficients $\ge 2$.
It is easily seen that the coefficients $b_{j}$ of $p_{V}^{k}$ then satisfy $b_{j} > k$ for $k \in \N$. This implies that for $k > \frac{n}{2}$ the unitary groups $U(b_{j})$ will all have dimension $b_{j} > \frac{n}{2}$. Now let $U_{\infty}$ be the colimit over the inclusions $U(n) \hookrightarrow U(n+1)$ that add a $1$ in the lower right corner. In the above situation $n$ falls into the stable range for the homotopy groups, which implies that $\pi_{n}(U(b_{j})) \cong \pi_{n}(U_{\infty})$, which is isomorphic to $\Z$ if $n$ is odd and vanishes if $n$ is even. 

Therefore for $k > \frac{n}{2}$ each term in the sequence which computes the group $\pi_{n}(U(D^{\Z/p\Z}))$ is isomorphic to the direct sum of one copy of $\Z$ for every irreducible $\Z/p\Z$-representation, i.e., is isomorphic to the representation ring $R(\Z/p\Z)$ as $R(\Z/p\Z)$-modules. Since the connecting homomorphism is given by multiplication by $p_{V}$, the result follows after taking the direct limit. 
\end{proof}

\begin{remark}
\label{K1}
The proof of Theorem \ref{homotopyU} shows in particular that the natural map
\[\pi_{n}(U(D^{\Z/p\Z})) \rightarrow K_{1}(C(S^n,D^{\Z/p\Z}))\]
given by mapping the class of $\gamma: S^{n} \rightarrow U(D^{\Z/p\Z})$ to the class of the corresponding unitary $u_{\gamma} \in U(C(S^n,D^{\Z/p\Z}))$ is an isomorphism. 
\end{remark}

\begin{corollary} \label{cor:pikAut}
    We have 
    \[
        \pi_n(\eqAut{\Z/p\Z}{D \otimes \bK}) \cong \begin{cases}
            GL_1(R(\Z/p\Z)[p_V^{-1}])_+ & \text{if } n = 0, \\
            R(\Z/p\Z)[p_V^{-1}] &\text{if } n > 0 \text{ and } n \text{ is even},\\
            0 & \text{if } n \text{ is odd}. 
        \end{cases}
    \]
\end{corollary}

\begin{proof}
This follows by combining Lemma~\ref{units+}, Proposition~\ref{ProjBU}, Proposition~\ref{AutProj}, and Theorem~\ref{homotopyU}.
\end{proof}

\section{$\Z/p\Z$-equivariant $\Aut{D \otimes \bK}$-bundles and units of $KU^D$}
\label{final_section}

\subsection{The classifying space of equivariant $\Aut{D \otimes \bK}$-bundles}
\label{final_1}
In this section we will take a closer look at equivariant bundles of $C^*$-algebras with fibre given by the $G$-$C^*$-algebra $D \otimes \bK$. We will assume that they are locally trivial in the following sense:
\begin{definition}
    Let $G$ be a finite group. Let $A$ be a $G$-$C^*$-algebra and denote the action by $\alpha$. Let $X$ be a topological space and let $\pi \colon \cA \to X$ be a locally trivial $C^*$-algebra bundle with fibre $A$ over $X$. We will say that $\cA$ is a \emph{locally trivial $(G,A,\alpha)$-bundle} (or a \emph{locally trivial $(G,A)$-bundle} if the action is clear) if 
    \begin{enumerate}[a)]
        \item $G$ acts from the left on the total space $\cA$ in such a way that $\pi$ is $G$-equivariant,
        \item for every $x \in X$ there exists a $G$-invariant open neighbourhood $U \subseteq X$ and a $G$-equivariant homeomorphism 
        \[
            \varphi_U \colon U \times A \to \left.\cA\right|_U\ ,
        \]
        where $\left.\cA\right|_U = \pi^{-1}(U)$ and $G$ acts on $U \times A$ via $g \cdot (x,a) = (g \cdot x, \alpha_g(a))$.
    \end{enumerate}
\end{definition}

To save space we will drop the adjective ``locally trivial'' from the notation in the following and instead mention the cases explicitly where we do not assume it. By elementary bundle theory each $C^*$-algebra bundle $\cA \to X$ with fibre $A$ gives rise to a principal $\Aut{A}$-bundle $\cP \to X$ such that there is a bundle isomorphism
\(
    \cA \cong \cP \times_{\Aut{A}} A
\). In fact, this construction induces a bijection between isomorphism classes of $C^*$-algebra bundles with fibre $A$ and principal $\Aut{A}$-bundles. 

We outline a few details of the construction of $\cP$: Let $\cA_x = \pi^{-1}(\{x\})$. The fibre $\cP_x$ of $\cP$ over $x \in X$ is given by
\begin{equation} \label{eqn:Px}
    \cP_x = \text{Iso}(A, \cA_x)\ ,
\end{equation}
which has a canonical right $\Aut{A}$-action, and the topology on $\cP$ is fixed by the property that a local trivialisation $\varphi_U \colon U \times A \to \left.\cA\right|_U$ over $U$ gives rise to a homeomorphism
\[
    U \times \Aut{A} \to \left.\cP\right|_U  \ .
\]
Now assume that $\cA \to X$ is a $(G,A)$-bundle. We will see that the associated principal $\Aut{A}$-bundle is equivariant in the following sense:
\begin{definition} \label{def:GAut-bundle}
    Let $G$ be a finite group, let $X$ be a topological space and let $A$ be a $G$-$C^*$-algebra. A $(G,\Aut{A})$-bundle is a locally trivial principal $\Aut{A}$-bundle $\cP \to X$ together with a left $G$-action by bundle maps such that each $x \in X$ has a $G$-invariant trivialising neighbourhood $U \subseteq X$ and a $G$-equivariant trivialisation
    \[
        \left.\cP\right|_U \to U \times \Aut{A}
    \]
    where $G$ acts from the left on $\Aut{A}$ via the action on $A$ and diagonally on $U \times \Aut{A}$.
\end{definition}

Note that $(G,\Aut{A}$)-bundles are locally trivial equivariant bundles in the sense of \cite[Def.~2.3]{paper:tomDieck} for the local object $\Aut{A} \to \ast$.

\begin{lemma}
    Let $G$ be a finite group and let $A$ be a $G$-$C^*$-algebra. The principal $\Aut{A}$-bundle associated to a $(G,A)$-bundle is a $(G,\Aut{A})$-bundle. Conversely, given a $(G,\Aut{A})$-bundle $\cP \to X$, then the associated $C^*$-algebra bundle is a $(G,A)$-bundle over $X$. 
\end{lemma}

\begin{proof}
The action of $G$ on $\cA$ induces a corresponding action on $\cP$ by post-composition (see \eqref{eqn:Px}) such that $G$ acts by bundle maps. Fix $x \in X$. Choose a $G$-invariant trivialising neighbourhood $U \subseteq X$ and a $G$-equivariant homeomorphism 
\[
    \varphi_U \colon U \times A \to \left.\cA\right|_U\ .
\]
The map $\varphi_U$ induces a $G$-equivariant trivialisation
\[
    \psi_U \colon U \times \Aut{A} \to \left.\cP\right|_U \quad , \quad (y, \beta) \mapsto \varphi_{U,y} \circ \beta\ .
\]
Conversely, assume that $\cP \to X$ is a $(G,\Aut{A})$-bundle. A $G$-equivariant trivialisation 
$\psi_U \colon U \times \Aut{A} \to \left.\cP\right|_U$ gives rise to a corresponding one for $\cA = \cP \times_{\Aut{A}} A$ defined by 
\[
    \varphi_U \colon U \times A \to \left.\cA\right|_U  \quad , \quad (y, a) \mapsto [\psi_U(x,\id{}),a] \ . \qedhere
\]
\end{proof} 

Following \cite{paper:tomDieck, book:MayClass} we will now construct a universal $(G,\Aut{A})$-bundle  $E\Aut{A} \to B\Aut{A}$ over the classifying space $B\Aut{A}$. Let
\[
    \cEG{n} = \Aut{A}^{n+1} \quad \text{ and } \quad \cBG{n} = \Aut{A}^n
\]
The face maps $d_i \colon \cEG{n} \to \cEG{n-1}$ are given by 
\[
    d_i(\beta_0, \dots, \beta_{n}) = (\beta_0, \dots, \beta_{i-1} \circ \beta_i, \dots,  \beta_{n}) \quad \text{ for } i > 0
\] 
and $d_0(\beta_0, \dots, \beta_{n}) = (\beta_1, \dots, \beta_{n})$. The degeneracy maps $s_i \colon \cEG{n} \to \cEG{n+1}$ are 
\[
    s_i(\beta_0, \dots, \beta_{n}) = (\beta_0, \dots, \beta_{i-1}, \id{}, \beta_i, \dots, \beta_{n})\ .
\]
The face and degeneracy maps for $\cBG{n}$ are defined similarly except for $d_n$, which drops the last component instead of composing. Both of these are simplicial $G$-spaces with respect to the action that is given by conjugation on the first $n$ factors of $\cEG{n}$ and $\cBG{n}$ and by the left action on the last factor of $\cEG{n}$. In the case of $\cEG{n}$ this is 
\[
    g \cdot (\beta_0, \dots, \beta_n) = (\alpha_g\beta_0\alpha_g^{-1}, \dots, \alpha_g\beta_{n-1}\alpha_g^{-1}, \alpha_g\beta_n)\ .
\]
Note that this action is compatible with the structure maps and that the obvious projection map $\cEG{\bullet} \to \cBG{\bullet}$ onto the first $n$ factors is $G$-equivariant. 

The two simplicial spaces $\cEG{\bullet}$ and $\cBG{\bullet}$ agree with the bar constructions $B(\ast, \Aut{A}_c, \Aut{A}_l)$ and $B(\ast, \Aut{A}_c, \ast)$ in the category of $G$-spaces, respectively, where $\Aut{A}_c$ is $\Aut{A}$ equipped with the conjugation action and $\Aut{A}_l$ is the same space equipped with the left action. Alternatively, the space $B\Aut{A}$ can be also viewed as the geometric realisation of the topological category associated to the group $\Aut{A}$. Let 
\[
    E\Aut{A} = \lvert \cEG{\bullet} \rvert \quad , \quad
    B\Aut{A} = \lvert \cBG{\bullet} \rvert\ .
\]
The projection map $\cEG{\bullet} \to \cBG{\bullet}$ induces a continuous $G$-equivariant map $\pi \colon E\Aut{A} \to B\Aut{A}$ on the geometric realisation.

\begin{prop}
    Let $G = \Z/p\Z$. Let $V$ be a finite-di\-men\-sional unitary $G$-representation. Let $D = \Endo{V}^{\otimes \infty}$ and $\bK = \bK(\ell^2(G) \otimes H_0)$. Then 
    \[
        \pi \colon E\Aut{D \otimes \bK} \to B\Aut{D \otimes \bK}
    \]
    is a universal $(G,\Aut{D \otimes \bK})$-bundle. 
\end{prop}

\begin{proof}
The inclusion $\{\id{}\} \to \Aut{D \otimes \bK}$ is a $G$-equivariant cofibration by Lemma~\ref{lem:G-cofib}. By \cite[Thm.~8.2]{book:MayClass} the map $\pi \colon E\Aut{D \otimes \bK} \to B\Aut{D \otimes \bK}$ is thus a principal $\Aut{D \otimes \bK}$-bundle. The construction of the local trivialisations in \cite[Thm.~8.2]{book:MayClass} is based on the observation that the $n$th simplicial filtration step included in the full space $E\Aut{D \otimes \bK}$ is a cofibration (see also \cite[Thm.~7.6]{book:MayClass}). Since this is true $G$-equivariantly here, the trivialisation described at the end of the proof of \cite[Thm.~8.2]{book:MayClass} is in fact $G$-equivariant in the sense of Definition ~\ref{def:GAut-bundle}. Therefore $E\Aut{D\otimes \bK} \to B\Aut{D \otimes \bK}$ is a $(G,\Aut{D \otimes \bK})$-bundle. 

To see that it is universal we may follow \cite[Thm.~5.1]{paper:tomDieck}: Recall that a local object $X \to G/H$ is a $(G,\Aut{A})$-bundle over the $G$-space~$G/H$ for a subgroup $H \subseteq G$. We have to show that for every local object $X \to G/H$ featuring in our local triviality condition the associated bundle 
\begin{equation} \label{eqn:EAutX}
    (E\Aut{D \otimes \bK} \times X)/\Aut{D \otimes \bK} \to G/H
\end{equation}
is $G$-shrinkable. In our case there is only one local object to check, namely $X = \Aut{D \otimes \bK} \to \ast = G/G$, i.e.\ the right $\Aut{D\otimes \bK}$-bundle over the point with the given left action of $G$ on $\Aut{D \otimes \bK}$. In this case the bundle in \eqref{eqn:EAutX} is homeomorphic to $E\Aut{D \otimes \bK} \to \ast$ with the $G$-action given by conjugation, i.e.\ the $G$-space given by $B(\ast,\Aut{D \otimes \bK}_c, \Aut{D \otimes \bK}_c)$. Let $H \subseteq G$ be a subgroup. Taking fixed points is a finite limit and therefore commutes with geometric realisation. Combining these observations we have
\[
    E\Aut{D \otimes \bK}^H \cong E\eqAut{H}{D \otimes \bK}\ ,
\]
which is contractible. Hence, $E\Aut{D \otimes \bK}$ is weakly $G$-contractible and the result follows from Lemma ~\ref{lem:GCW-complex} and \cite[Prop.~13.2]{paper:Murayama}.
%The bundle \eqref{eqn:EAutX} in case \eqref{it:p1} is $G$-shrinkable if and only if $E\Aut{D \otimes \bK}$ is contractible, which follows just as in the non-equivariant setting (see for example \cite[Prop.~7.5]{book:MayClass}). 
\end{proof}

\begin{remark}
    The proof of universality shows that $E\Aut{D \otimes \bK}$ is in fact a model for $E_\mathcal{F(R)}\Aut{D \otimes \bK}$ for the family $\mathcal{R}$ of local representations in the sense of \cite[Def.~3.4]{paper:LueckUribe} that consists of all conjugates of the action $\alpha \colon G \to \Aut{D \otimes \bK}$ and the trivial homomorphism $\{\id{}\} \to \Aut{D \otimes \bK}$ (see \cite[Thm.~11.5]{paper:LueckUribe}). Note, however, that the equivariant principal bundles considered in \cite{paper:LueckUribe} are less restrictive than our definition. In particular, \cite[Def.~2.1]{paper:LueckUribe} only demands that local triviality holds non-equivariantly, whereas Definition~\ref{def:GAut-bundle} asks for $G$-equivariant local trivialisations (as in \cite{paper:tomDieck}). With the more flexible notion, Condition~(H) (see \cite[Def.~6.1]{paper:LueckUribe}) is needed to ensure homotopy invariance. It is not clear to us, if Condition~(H) holds in our case.
\end{remark}

\noindent Note that the $G$-space $\Aut{D \otimes \bK}$ has two equivariant deloopings: 
\begin{enumerate}
    \item The classifying space $B\Aut{D \otimes \bK}$ equipped with the $G$-action induced by the conjugation action of $G$ on $\Aut{D \otimes \bK}$ defined above. 
    \item The $\Gamma$-$G$-space delooping given by 
    \[
        B_\otimes\Aut{D\otimes \bK} = (EH_\cI\cG_D)_1
    \]
    induced by the tensor product structure on automorphisms.
\end{enumerate}

\begin{lemma} \label{lem:comp_vs_tensor}
Let $G=\Z/p\Z$ and let $D$ and $\bK$ be the $G$-$C^*$-algebras defined above. Then there is a weak $G$-equivalence
\[
    B\Aut{D\otimes \bK} \simeq B_\otimes\Aut{D \otimes \bK}\ .
\]
As a consequence, $B_\otimes\Aut{D \otimes \bK}$ is a classifying space for $G$-equivariant $D \otimes \bK$-bundles over finite CW-complexes.
\end{lemma}

\begin{proof}
    The proof is an equivariant version of \cite[Thm.~3.6]{paper:DP-Units} and can be reduced to that statement: With $\mu = \otimes$ induced by the tensor product and $\nu = \circ$ given by composition there is an intermediate space $B_\otimes B\Aut{D \otimes \bK}$ constructed in the proof of \cite[Thm.~3.6]{paper:DP-Units}. We can carry out this construction in the category of $G$-spaces. Then it suffices to check that for a subgroup $H \subseteq G$ the induced maps on fixed points
    \[
        \begin{tikzcd}
        B_\otimes\Aut{D\otimes \bK}^H \ar[r] & 
        \Omega B_\otimes B\Aut{D \otimes \bK}^H &
        B\Aut{D \otimes \bK}^H \ar[l]
    \end{tikzcd}    
    \]
    is a weak equivalence. But because fixed points commute with limits and the group action is by conjugation the above boils down to the following sequence of maps
    \[
    \begin{tikzcd}
        B_\otimes\!\eqAut{H}{D\otimes \bK} \ar[r,"\simeq"] & 
        \Omega B_\otimes B\!\eqAut{H}{D \otimes \bK} &
        B\!\eqAut{H}{D \otimes \bK} \ar[l,"\simeq" above]
    \end{tikzcd}\ ,
    \]
    all of which are weak equivalences by the original \cite[Thm.~3.6]{paper:DP-Units}.
\end{proof}

\subsection{Comparison between $\cG_D$ and $\textnormal{GL}^{\bullet}_1 KU^D$}
\label{final_2}

Let $e \in \bK$ be a projection of rank one fixed by $G$. We can naturally define a collection of maps
\[\theta^H_n\colon \eqAut{H}{(D \otimes \bK)^{\otimes n}} \rightarrow (\Omega^n KU^{D}_{n})^H=\!\hom_{\text{gr}}(\grS, C_{0}(\R^n) \otimes \C \ell_1 \otimes (D \otimes \bK)^H)\]
\[ \alpha \mapsto (f \mapsto \hat{\eta}_{n}(f) \otimes \alpha(1 \otimes e))\]
for all $H \subseteq G$. Since $\cG_D$ has compatible inverses (see Lemma \ref{GD_stable}), these maps factor over a morphism
\[\theta: \cG_D \rightarrow \text{GL}_1^{\bullet} KU^D\]
of commutative $\cI$-$O_G$-monoids. Corollary~\ref{cor:pikAut} suggests it might be worth to investigate whether $\theta$ is a weak equivalence when $G=\Z/p\Z$. This turns out to be true if one takes the order structure on $\pi_0(KU^D) \cong K_0(D)$ into account. Assuming that $\pi_0(R)$ has an order structure we define $GL_1^\bullet(R)_+$ by a pullback diagram analogous to the one in \eqref{eqn:GL1_diag} with $GL_1(\pi_0(R))$ replaced by $GL_1(\pi_0(R))_+$. Let $gl_1(R)_+ = EH_\cI GL_1^\bullet(R)_+$ be the associated $G$-spectrum.

%Let us denote by $gl_1(R)_+$ the $G$-spectrum obtained as in Section \ref{equnits}, when we assume that $\pi_0 R$ has an order structure and we restrict to the positive components $(\pi_0 R)_+$. 
%It turns out this is almost the case, and that the equivalence can be lifted to the corresponding $\Z/p\Z$-spectra.

\begin{theorem}\label{thm:main_thm}
There is a map of $\Z/p\Z$-spectra
\[ EH_{\cI}\Aut{D \otimes \bK} \rightarrow gl_{1}(KU^{D})\]
which is an isomorphism on all higher equivariant homotopy groups $\pi^{H}_{n}$ with $n >0$, and the inclusion $\text{GL}_1(K_{0}^{H}(D))_{+} \hookrightarrow \text{GL}_1(K_{0}^{H}(D))$ on $\pi_{0}^{H}$, for $H \subseteq \Z/p\Z$. In particular, we have an equivalence of $\Z/p\Z$-spectra
\[
EH_{\cI}\Aut{D \otimes \bK} \simeq gl_{1}(KU^{D})_+\ .
\]
\end{theorem}

\begin{proof}
By Lemma \ref{deloop} it is enough to check that the maps
\[ \theta_1^H: \eqAut{H}{D \otimes \bK} \rightarrow (\Omega KU_{1}^{D})^{H}\]
defined above enjoy the prescribed properties, for $H \in \{\{e\}, \Z/p\Z\}$.
Observe that $\theta_1^H$ factors through
% https://q.uiver.app/#q=WzAsMyxbMCwwLCJcXHRleHR7QXV0fSgoRCBcXG90aW1lcyBcXG1hdGhiYntLfSleSCkiXSxbMiwwLCJcXHRleHR7aG9tfV97XFx0ZXh0e2dyfX0oXFxoYXR7U30sIENfezB9KFxcUikgXFxvdGltZXMgXFxtYXRoYmJ7Q30gXFxlbGxfMSBcXG90aW1lcyAoRCBcXG90aW1lcyBcXG1hdGhiYntLfSleSCkiXSxbMCwxLCJcXHRleHR7UHJvan1fezEgXFxvdGltZXMgZX0oKEQgXFxvdGltZXMgXFxtYXRoYmJ7S30pXkgpIl0sWzAsMSwiXFx0aGV0YSJdLFswLDIsIlxcUGhpIiwyXSxbMiwxLCJcXFBzaSIsMl1d
\[\begin{tikzcd}
	\eqAut{H}{D \otimes \bK} && {\text{hom}_{\text{gr}}(\widehat{\mathcal{S}}, C_{0}(\R) \otimes \mathbb{C} \ell_1 \otimes (D \otimes \mathbb{K})^H)} \\
	{\text{Proj}((D \otimes \mathbb{K})^H)}
	\arrow["\theta_1^H", from=1-1, to=1-3]
	\arrow["\Phi"', from=1-1, to=2-1]
	\arrow["\Psi"', from=2-1, to=1-3]
\end{tikzcd}\]
where $\Psi(p)= (f \mapsto \eta_{1}(f) \otimes p)$. By Proposition \ref{AutProj} the map $\Phi$ is a homotopy equivalence when the domain is restricted to $\eqAutId{H}{D \otimes \bK}$ and the codomain to the component of $1 \otimes e$. Recall that $\hat{\eta}_1 \in \hom_{\text{gr}}(\grS, C_0(\R, \Cln{1}))$ is the Bott element. Note that $\Psi$ factors in turn as
\[
\begin{tikzcd}[column sep=0.5cm,font=\footnotesize]
\Proj{}{(D \otimes \bK)^H} \ar[r] & \hom_{\text{gr}}(\grS, (D \otimes \bK)^H) \ar[r] & \hom_{\text{gr}}(\grS, C_{0}(\R, \C \ell_1) \otimes (D \otimes \bK)^H).
\end{tikzcd}
\]

Here the first map sends a projection $p$ to $\epsilon \cdot
\left(\begin{smallmatrix}
p & 0\\
0 & 0
\end{smallmatrix}\right)$. The second map sends $\varphi$ to $\varphi \otimes \hat{\eta_1} \circ \Delta$ and shifts the grading to $C_{0}(\R) \otimes \C \ell_1$, and it is an isomorphism on $\pi_0$ (since it induces multiplication by the Bott element). Therefore the discussion after \cite[Thm. 4.7]{paper:Trout}, plus Lemma~\ref{units+} show that $\theta_1^H$ induces the inclusion $\text{GL}_1(K_{0}^{H}(D))_{+} \hookrightarrow \text{GL}_1(K_{0}^{H}(D))$ on $\pi_{0}$. 

To check $\pi_n$ for $n > 0$ we may restrict $\Phi$ to the equivalence
\[
    \Phi \colon \eqAutId{H}{D \otimes \bK} \to \Proj{1 \otimes e}{(D \otimes \bK)^H}
\]
and $\Psi$ accordingly. Consider the following commutative square
% https://q.uiver.app/#q=WzAsNCxbMCwwLCJcXHBpX3tufVxcdGV4dHtQcm9qfSgoRCBcXG90aW1lcyBcXG1hdGhiYntLfSleSCkiXSxbMiwwLCJcXHBpX3tufVxcT21lZ2EgS1Vee0R9X3sxfSJdLFswLDEsIktfMChDXzAoU15uLCopIFxcb3RpbWVzIEReSCkiXSxbMiwxLCJLJyhDX3swfShcXG1hdGhiYntSfSkgXFxvdGltZXMgXFxtYXRoYmJ7Q30gXFxlbGwgX3sxfSBcXG90aW1lcyBDKFNee259KSBcXG90aW1lcyBEXkgpIl0sWzAsMSwiXFxwc2lfKiJdLFswLDIsIlxcYmV0YSIsMl0sWzIsMywiXFxjb25nIiwyXSxbMSwzLCJcXGNvbmciXV0=
\[\begin{tikzcd}
	{\pi_{n}(\text{Proj}_{1 \otimes e}((D \otimes \mathbb{K})^H))} && {\pi_{n}((\Omega KU^{D}_{1})^H)} \\
	{K_0(C_0(S^n,*) \otimes D^H)} && {K'(C_{0}(\mathbb{R}) \otimes \mathbb{C} \ell _{1} \otimes C(S^{n}) \otimes D^H)}
	\arrow["{\Psi_*}", from=1-1, to=1-3]
	\arrow["\beta"', from=1-1, to=2-1]
	\arrow["\cong"', from=2-1, to=2-3]
	\arrow["\cong", from=1-3, to=2-3]
\end{tikzcd}\]
where we are using the notation $K'(A)= \pi_{0}(\hom_{\text{gr}}(\grS,A \otimes \bK))$ introduced in \cite{paper:Trout}. Note that any element $\gamma: S^n \rightarrow \text{Proj}_{1 \otimes e}((D \otimes \mathbb{K})^H)$ induces a projection $p_{\gamma} \in  C(S^{n}) \otimes D^H$, and set $\beta(\gamma)=[p_{\gamma}] - [1_{C(S^n)} \otimes 1 \otimes e]$. The bottom horizontal map sends $[p]-[q]$ to $\left[\left(f \mapsto \hat{\eta}_1 \otimes
\left(\begin{smallmatrix}
p & 0\\
0 & q
\end{smallmatrix}\right)\right)\right]$ and is an isomorphism by Bott periodicity. Finally, note that any element in $\pi_n ((\Omega KU_1^D)^H)$ defines a map $\varphi \in \hom_{\text{gr}}(\grS,C_0(\R) \otimes \C \ell_1 \otimes C(S^n) \otimes (D \otimes \bK)^H)$. The right hand vertical arrow maps $[\varphi]$ to 
$\left[\left(
\begin{smallmatrix}
\varphi & 0\\
0 & 1_{C(S^n)} \otimes \hat{\eta}_1
\end{smallmatrix}\right)\right]$ and has an inverse given by $\psi \mapsto \psi \oplus 
\left(\begin{smallmatrix}
1_{C(S^n)} \otimes \hat{\eta}_1& 0\\
0 & 0
\end{smallmatrix}\right)$, where $\oplus$ is the addition operation described in~\cite{paper:Trout}. Note that this is a basepoint correction, because in $\Psi_*$ we use a basepoint for $\pi_n$ in the $[1 \otimes e]$-component on both sides instead of the usual $0$-component.

The last step is to prove that $\beta$ is an isomorphism as well. It will follow that $\Psi$ induces isomorphisms on all higher homotopy groups, hence so does~$\theta$. For this purpose, let us consider the following square:
% https://q.uiver.app/#q=WzAsNCxbMCwwLCJcXHBpX3tuLTF9KFUoRF5IKSkiXSxbMCwxLCJLXzEoQ18wKFNee24tMX0sKikgXFxvdGltZXMgRF5IKSJdLFsxLDAsIlxccGlfbihcXHRleHR7UHJvan1fezEgXFxvdGltZXMgZX0oKEQgXFxvdGltZXMgXFxtYXRoYmJ7S30pXkgpKSJdLFsxLDEsIktfMChDXzAoU15uLCopIFxcb3RpbWVzIEReSCkiXSxbMiwzLCJcXGJldGEiXSxbMywxLCIiLDIseyJzdHlsZSI6eyJib2R5Ijp7Im5hbWUiOiJub25lIn0sImhlYWQiOnsibmFtZSI6Im5vbmUifX19XSxbMCwxLCJcXGNvbmciLDJdLFswLDIsIlxcY29uZyJdLFsxLDMsIlxcY29uZyIsMl1d
\[\begin{tikzcd}
	{\pi_{n-1}(U(D^H))} & {\pi_n(\text{Proj}_{1 \otimes e}((D \otimes \mathbb{K})^H))} \\
	{K_1(C_0(S^{n-1},*) \otimes D^H)} & {K_0(C_0(S^n,*) \otimes D^H)}
	\arrow["\beta", from=1-2, to=2-2]
	\arrow["\cong"', from=1-1, to=2-1]
	\arrow["\cong", from=1-1, to=1-2]
	\arrow["\cong"', from=2-1, to=2-2]
\end{tikzcd}\]
Here the top horizontal arrow is the inverse of the isomorphism described in Remark \ref{RemProjBU} composed with the canonical isomorphism $\pi_n(BG) \cong \pi_{n-1}(G)$. This map is defined by constructing a principal $G$-bundle over $S^n$ with the prescribed transition map $\psi \in \pi_{n-1}(G)$, explicitly given by $P_{\psi}=S^{n-1} \times I \times G$ modulo the equivalence relation $\sim$ generated by
\begin{equation} \label{eqn:eqrel}
    (z,1,g) \sim (z,0,\psi(z)\cdot g) \text{ and }(z_0,t,g) \sim (z_0,0,g)\ ,
\end{equation}
with the quotient map to $\Sigma S^{n-1} \cong S^n$ as projection map (see \cite[Sec.~4.4]{book:Naber}). Recall $\beta$ maps $[\gamma]$ to $[p_{\gamma}] - [1_{C(S^n)} \otimes 1 \otimes e]$ as in the diagram above. We can interpret the composition of these two arrows as the map sending $\gamma \in \pi_{n-1}(U(D^H))$ to the formal difference $[\mathcal{H}_{\gamma}]-[\underline{\mathcal{H}}]$ of right Hilbert $D^H$-module bundles over $S^n$, where
\[ \mathcal{H}_{\gamma}=\bigl(S^{n-1} \times I \times U(D^H)\bigr) \,/ \sim \ ,\]
and the equivalence relation is as in \eqref{eqn:eqrel} with $\psi=\gamma$. On the other hand, if we start from the left hand vertical arrow we encounter the isomorphism described in Remark \ref{K1}, and the bottom horizontal arrow is the suspension isomorphism in $K$-theory (see Example \ref{indexmap}). It is easy to check that this composite also maps $\gamma$ to $[\mathcal{H}_{\gamma}]-[\underline{\mathcal{H}}]$. Therefore the diagram commutes, implying that the map $\beta$ is an isomorphism, thus concluding the proof.
\end{proof}

Our final result combines Lemma~\ref{lem:comp_vs_tensor} with Theorem~\ref{thm:main_thm}.

\begin{corollary} \label{cor:bdl_classification}
    $\Aut{D \otimes \bK}$ is an equivariant infinite loop space with associated $\Z/p\Z$-equivariant cohomology theory $E_D^*(X)=gl_1(KU^D)_+^*(X)$, and
    \[E^0_D(X)=[X,\Aut{D \otimes \bK}]^{\Z/p\Z} \quad \text{and} \quad E^1_D(X)=[X,B\Aut{D \otimes \bK}]^{\Z/p\Z}.\]
    In particular, isomorphism classes of $\Z/p\Z$-equivariant $C^*$-algebra bundles with fibres isomorphic to the $\Z/p\Z$-algebra $D \otimes \bK$ over the finite $\Z/p\Z$-$CW$-complex $X$ form a group with respect to the fibrewise tensor product that is isomorphic to $E^1_D(X)$.
\end{corollary}

\appendix
\section{The equivariant homotopy type of $\Aut{D \otimes \bK}$}
Let $G = \Z/p\Z$, denote by $D = \Endo{V}^{\otimes \infty}$ the UHF-algebra associated to a finite-dimensional unitary $G$-representation $V$ and let $\bK = \bK(\ell^2(G) \otimes H_0)$ for an infinite-dimensional separable Hilbert space $H_0$. Let $\Proj{}{D \otimes \bK}$ be the space of projections in $D \otimes \bK$ equipped with the subspace topology and denote by $\Aut{D \otimes \bK}$ the group of $*$-automorphisms with the point-norm topology considered as a $G$-space with the induced action.

In this appendix we will show that the automorphism group $\Aut{D \otimes \bK}$ is equivariantly well-pointed in the sense that the inclusion of the identity is an equivariant cofibration. Apart from this, we will also prove a few technical results that might be of independent interest: In particular, we will see that $\Aut{D \otimes \bK}$ is $G$-homotopy equivalent to a $G$-CW-complex. 

\begin{lemma} \label{lem:GCW-complex}
Let $G$, $V$, $D$ and $\bK$ be as above. The topological group $\Aut{D \otimes \bK}$ equipped with the conjugation action of $G$ has the $G$-homotopy type of a $G$-CW-complex.
\end{lemma}

\begin{proof}
Pick a projection $e \in \bK$ whose image is the one-dimensional trivial $G$-representation. Let $q_0 = 1 \otimes e \in D \otimes \bK$. Let $P_D \colon O_G^{\op} \to \Top$ be the $O_G$-space defined by 
\[
    P_D(G/H) = \Proj{I_H}{(D \otimes \bK)^H}\ ,
\]
where $I_H = GL_1(K_0^H(D)_+)$. Let $A_D \colon O_G^\op \to \Top$ be the $O_G$-space associated to the $G$-space $\Aut{D\otimes \bK}$ equipped with the conjugation action, i.e.\ 
\[
    A_D(G/H) = (\Aut{D \otimes \bK})^H = \eqAut{H}{D \otimes \bK}\ .
\]
Let $\pi \colon A_D \to P_D$ be the map of $O_G$-spaces given by $\pi(\beta) = \beta(q_0)$ and note that $\pi$ provides a homotopy equivalence $A_D(G/H) \to P_D(G/H)$ for each subgroup $H \subseteq G$. For $H = \{e\}$ this is proven in \cite[Lem.~2.16 and Thm.~2.5]{paper:DP-Dixmier}, for $H = \Z/p\Z$ the proofs of \cite[Prop.~3.1, Lem.~2.9 and Thm.~2.6]{paper:EP-Circle} for circle actions carry over to the case of $\Z/p\Z$ . 

Let $\Phi \colon \text{Fun}(O_G^\op, \Top) \to \GTop$ be the homotopy inverse of the fixed-point functor defined in (\ref{eqn:elmendorf}). Note that $\Phi(X)$ can be written as the geometric realisation of a simplicial space obtained from the bar construction. The fact that it provides a homotopy inverse is witnessed by a simplicial homotopy equivalence. If we therefore use the fat geometric realisation (i.e.\ ignoring the degeneracy maps), the resulting $\Phi$ is still a homotopy inverse. Moreover, $\pi$ provides a level-wise $G$-homotopy equivalence of simplicial spaces
\[
    B_\bullet(A_D, O_G, M) \to B_\bullet(P_D, O_G, M)\ .
\]
The fat geometric realisation turns this into a $G$-homotopy equivalence. This is well-known in the non-equivariant setting (see \cite[Prop.~A.1~(ii)]{paper:SegalCatAndCoh}) and can be deduced from the fact that pushouts along cofibrations preserve homotopy equivalences and that colimits over cofibrations are homotopy invariant. Both of these statements are still true in the $G$-equivariant setting (see \cite[Thm.~1.1 and Thm.~1.2]{paper:Waner}). Each of the spaces $P_D(G/H)$ has the homotopy type of a CW-complex by \cite[Lem.~2.7]{paper:DP-Dixmier}. Therefore $\Phi(P_D)$ has the homotopy type of a $G$-CW-complex by \cite[Prop.~13.2]{paper:Murayama} and we have the following sequence of $G$-homotopy equivalences
\[
    \begin{tikzcd}
        \Aut{D \otimes \bK} & \ar[l,"\simeq" above] \Phi(A_D) \ar[r,"\simeq"] & \Phi(P_D)  \ .   
    \end{tikzcd}\qedhere
\]
\end{proof}

As we will see in the next lemma with respect to finite group actions the space of projections $\Proj{}{A}$ of a $G$-$C^*$-algebra also has the $G$-homotopy type of a $G$-CW-complex.
\begin{lemma}
    Let $G$ be a finite group and let $A$ be a $G$-$C^*$-algebra. The subspace $\Proj{}{A} \subset A$ has the $G$-homotopy type of a $G$-CW-complex. Let $I \subseteq \pi_0(\Proj{}{A})$ be a $G$-invariant subset and let $\Proj{I}{A}$ be the pullback of $\Proj{}{A}$ over $I$. Then the space
    \[
        \Proj{I}{A} = \coprod_{[q] \in I} \Proj{q}{A}
    \]
    has the $G$-homotopy type of a $G$-CW-complex as well.
\end{lemma}

\begin{proof}
    By \cite{paper:Antonyan} the real Banach space $A_{sa}$ of self-adjoint elements in the algebra $A$ is a $G$-ANE ($G$-equivariant absolute neighbourhood extensor). Therefore the $G$-invariant open subset 
    \[
        A_{sa}^0 = \{ x \in A_{sa} \ |\ \tfrac{1}{2} \notin \sigma(x) \} \subset A_{sa}
    \]
    is a $G$-ANE as well \cite[Prop.~6.1]{paper:Murayama}. By \cite[Prop.~6.4]{paper:Murayama} it is then also a $G$-ANR ($G$-equivariant absolute neighbourhood retract) and \cite[Thm.~13.3]{paper:Murayama} implies that it has the $G$-homotopy type of a $G$-CW-complex. Note that 
    \[
        \Proj{}{A} \subset A_{sa}^0\ .
    \]
    Applying continuous functional calculus with the map that is $0$ on $(-\infty, \tfrac{1}{2})$ and $1$ on $(\tfrac{1}{2}, \infty)$ to the elements of $A_{sa}^0$ we obtain a $G$-equivariant map 
    \[
        \pi \colon A_{sa}^0 \to \Proj{}{A}\ ,
    \]
    that restricts to the identity on $\Proj{}{A}$, In particular, \ $\Proj{}{A}$ is dominated by a space with the $G$-homotopy type of a $G$-CW-complex. Therefore \cite[Lem.~4.7]{paper:Waner} implies that it has the $G$-homotopy type of a $G$-CW-complex itself\footnote{The lemma in \cite{paper:Waner} is stated for spaces that are dominated by $G$-CW-complexes, but the argument actually only requires the dominating space to have the correct $G$-homotopy type.}. The second statement is a consequence of the first.
\end{proof}

In the case of finite group actions evaluating an automorphism at a fixed projection gives an equivariant Hurewicz fibration as we will see below.

\begin{lemma} \label{lem:Hurewicz}
Let $G$ be a finite group and let $A$ be a $G$-$C^*$-algebra. Let $q_0 \in \Proj{}{A}$ be a projection, which is fixed by $G$.
Let 
\[
    \Xi \colon \Aut{A} \to \Proj{I}{A} \quad, \quad \beta \mapsto \beta(q_0)
\]
where $I \subseteq \pi_0(\Proj{}{A})$ is the image of $\Xi_*$ on $\pi_0$. Then $\Xi$ is a Hurewicz $G$-fibration (where $G$ acts on the domain by conjugation). 
\end{lemma}

\begin{proof}
We will show that $\Xi \colon \Aut{A} \to \Proj{I}{A}$ is a locally trivial $G$-fibration in the sense of \cite[Sec.~5]{paper:BoothHoskins}. Since the codomain is a metric space, hence paracompact, the result follows from \cite[Lem.~5.1]{paper:BoothHoskins}, \cite[Lem.~1.12]{paper:LashofEquivBundles} and \cite[Lem.~3.2.4]{paper:WanerClassFib}. Let $p_0 \in \Proj{I}{A}$ and $H = \text{stab}_G(p_0)\subseteq G$ be the stabiliser of $p_0$. For two projections $s, t \in \Proj{I}{A}$ we define 
\[
    \Autf{s}{t}{A} = \{ \beta \in \Aut{A}\ |\ \beta(s) = t \} \subseteq \Aut{A}\ .
\]
The group $H$ acts on $\Autf{q_0}{p_0}{A}$ by conjugation and 
\[
    G \times_H \Autf{q_0}{p_0}{A} \cong \Autf{q_0}{Gp_0}{A} \to G/H
\]
with $\Autf{q_0}{Gp_0}{A} = \{\beta \in \Aut{A}\ |\ \beta(q_0) \in G \cdot p_0\}$ is $G$-equivariant and a $G$-model trivial fibration \cite[Sec.~5]{paper:BoothHoskins}. Define
\[    
    U_0 = \{ q \in \Proj{I}{A}\ | \ \lVert q-p_0 \rVert < 1\} \ .
\]
The subspace $U_0$ is $H$-invariant. As a metric space $\Proj{I}{A}$ is completely regular. By \cite[Thm.~1.7.19]{paper:Palais} there exists a slice $V \subseteq \Proj{I}{A}$ with $p_0 \in V$ and $V \cap U_0$ provides another slice containing $p_0$. Let 
\[
    U = G \cdot (V \cap U_0) \cong G \times_H (V \cap U_0)\ .
\]
By \cite[Prop.~II.3.3.4]{book:BlackadarOpAlg} there is a continuous map
\(
    u_{p_0} \colon U_0 \to U(M(A))
\)
with the properties 
\begin{enumerate}[i)]
    \item \label{it:transport} $u_{p_0}(q)\,p_0\,u_{p_0}(q)^* = q$,
    \item \label{it:equiv} $\alpha_h(u_{p_0}(q)) = u_{p_0}(\alpha_h(q))$ for all $h \in H$.
\end{enumerate}
Now consider the continuous map
\[
   \kappa \colon G \times_H ((V \cap U_0) \times \Autf{q_0}{p_0}{A}) \to \Xi^{-1}(U)
\]
given by $\kappa([g, (q,\beta)]) = \alpha_g \circ \Ad_{u_{p_0}(q)} \circ \beta \circ \alpha_g^{-1}$. Since
\[
    \alpha_{gh} \circ \Ad_{u_{p_0}(\alpha_h^{-1}(q))} \circ\ (\alpha_h^{-1} \circ \beta \circ \alpha_h) \circ \alpha_{gh}^{-1} = \alpha_g \circ \Ad_{u_{p_0}(q)} \circ \beta \circ \alpha_g^{-1}
\]
this is well-defined. It is also equivariant with respect to the left action of~$G$ on the domain and the conjugation action on the codomain. To define the inverse map, choose coset representatives $S = \{g_0 = e, g_1, \dots, g_n\} \subset G$ for~$G/H$. Observe that for $\beta \in \Xi^{-1}(U)$ the projection $\beta(q_0) \in G \cdot (V \cap U_0)$ can be written as $\alpha_{g_\beta}(q_\beta)$ for a unique $g_\beta \in S$ and a projection~$q_\beta$ with $\lVert q_\beta - p_0 \rVert < 1$. Define
\[
    \theta \colon \Xi^{-1}(U) \to G \times_H ((V \cap U_0) \times \Autf{q_0}{p_0}{A})
\]
by $\theta(\beta) = [g_\beta, (q_\beta, \Ad_{u_{p_0}(q_\beta)^*} \circ\ 
\alpha_{g_\beta}^{-1} \circ \beta \circ \alpha_{g_\beta})]$. A straightforward computation shows that $\theta \circ \kappa = \id{}$ and $\kappa \circ \theta = \id{}$, which implies that $\theta$ is $G$-equivariant as well. Hence, $\kappa$ is a local trivialisation.
\end{proof}

\begin{remark}
    The space $\Autf{q_0}{p_0}{A}$ is a right $\AutSt{q_0}{A}$-torsor. The left action of $G$ on $\Aut{A}$ commutes with the right action by $\AutSt{q_0}{A}$ and $\pi$ is equivariant with respect to those two actions. Replacing conjugation actions by left actions everywhere in the proof one can show that $\Xi \colon \Aut{A} \to \Proj{I}{A}$ is a $G$-equivariant principal bundle with structure group $\AutSt{q_0}{A}$ in analogy to \cite[Lem.~2.8]{paper:DP-Dixmier}. 
\end{remark}


\begin{lemma} \label{lem:G-cofib}
Let $G$, $V$, $D$ and $\bK$ be as in the first paragraph of the appendix. The topological group $\Aut{D \otimes \bK}$ equipped with the conjugation action of $G$ is equivariantly well-pointed in the sense that the inclusion $\{\id{}\} \to \Aut{D \otimes \bK}$ is a $G$-cofibration.
\end{lemma}

\begin{proof}
    By a straightforward $G$-equivariant generalisation of \cite[Thm.~2]{paper:Strom} (see also \cite[Lem.~2.3]{paper:Grant}) it suffices to construct a $G$-invariant continuous map $v \colon \AutId{D \otimes \bK} \to [0,1]$ such that 
    \begin{enumerate}
        \item $v^{-1}(0) = \{\id{}\}$,
        \item $U_v = v^{-1}([0,1)) = \{ \beta \in \AutId{D \otimes \bK} \ |\ v(\beta) < 1\}$ deformation retracts $G$-equivariantly to $\id{}$.
    \end{enumerate}
    Pick a countable dense subset $S = \{ a_k \in D\otimes \bK \ |\ k \in \N\}$, which is $G$-invariant as a set. The point-norm topology on $\Aut{D \otimes \bK}$ is the metric topology associated to the metric
    \[
        d(\beta_1, \beta_2) = \sum_{k \in \N} \frac{\lVert \beta_1(a_k) - \beta_2(a_k) \rVert}{2^k \lVert a_k \rVert}\ .
    \]
    Since each $\alpha_g$ gives a bijection $S \to S$ and $\lVert \alpha_g(b) \rVert = \lVert b \rVert$ for all $b \in D \otimes \bK$, this metric is $G$-invariant. Let $q_0 = 1 \otimes e \in D \otimes \bK$ be as in the proof of Lemma~\ref{lem:GCW-complex} and $\Xi$ be as in Lemma~\ref{lem:Hurewicz}. We may now proceed as in \cite[Prop.~2.26]{paper:DP-Dixmier}. Let
    \[
        v(\beta) = \max\{\min\{d(\alpha, \id{}), \tfrac{1}{2}\}, \min\{1, 2\lVert \beta(q_0) - q_0 \rVert\} \}\ ,
    \]
    which is $G$-invariant, continuous and satisfies $v^{-1}(0) = \{\id{}\}$. Note that
    \[
        U_v = \Xi^{-1}(W) \quad \text{with} \quad W = \{ q \in \Proj{q_0}{D \otimes \bK} \ |\ \lVert q - q_0 \rVert < \tfrac{1}{2} \}\ .
    \]
    As in the proof of Lemma~\ref{lem:Hurewicz} the map
    \[
        \kappa \colon W \times \AutSt{q_0}{D\otimes \bK} \to \Xi^{-1}(W)
    \]
    with $\kappa(q,\beta) = \Ad_{u_{q_0}(q)} \circ \beta$ is a $G$-equivariant homeomorphism. Hence, it suffices to prove that the domain equivariantly deformation retracts to $(q_0, \id{})$. Let $\chi$ be the characteristic function of $(\tfrac{1}{2},1]$. The $G$-equivariance of functional calculus and the choice of $W$ ensure that $h(q,t) = \chi((1-t)q + tq_0)$ is a well-defined equivariant retraction of $W$ to $q_0$. The homotopy described in \cite[Thm.~2.6]{paper:EP-Circle} uses the path \cite[eqn.~(9)]{paper:EP-Circle}
    \[
        \gamma \colon [0,1] \to \hom_G(D \otimes \bK, D \otimes \bK \otimes D \otimes \bK)
    \]
    connecting $a \otimes q_0$ to $q_0 \otimes a$ through isomorphisms in the interior of the interval. This carries over to $G=\Z/p\Z$ and shows that $\AutSt{q_0}{D \otimes \bK}$ equivariantly deformation retracts to $\id{}$. Combining both retractions gives the result.
\end{proof}

\bibliographystyle{plain}
\bibliography{main}

\end{document}
\begin{thebibliography}{10}

\bibitem{paper:Antonyan}
S.~A. Antonyan.
\newblock Equivariant generalization of {D}ugundji's theorem.
\newblock {\em Mat. Zametki}, 38(4):608--616, 636, 1985.

\bibitem{paper:AtiyahSegal}
Michael Atiyah and Graeme Segal.
\newblock Twisted {$K$}-theory.
\newblock {\em Ukr. Mat. Visn.}, 1(3):287--330, 2004.

\bibitem{book:BlackadarOpAlg}
B.~Blackadar.
\newblock {\em Operator algebras}, volume 122 of {\em Encyclopaedia of
  Mathematical Sciences}.
\newblock Springer-Verlag, Berlin, 2006.
\newblock Theory of $C^*$-algebras and von Neumann algebras, Operator Algebras
  and Non-commutative Geometry, III.

\bibitem{book:Blackadar}
Bruce Blackadar.
\newblock {\em {$K$}-theory for operator algebras}, volume~5 of {\em
  Mathematical Sciences Research Institute Publications}.
\newblock Cambridge University Press, Cambridge, second edition, 1998.

\bibitem{paper:BoothHoskins}
Peter~I. Booth and Philip~J. Hoskins.
\newblock On equivariant fibrations.
\newblock {\em Bull. Belg. Math. Soc. Simon Stevin}, 9(4):545--554, 2002.

\bibitem{paper:DP-Units}
Marius Dadarlat and Ulrich Pennig.
\newblock Unit spectra of {$K$}-theory from strongly self-absorbing
  {$C^*$}-algebras.
\newblock {\em Algebr. Geom. Topol.}, 15(1):137--168, 2015.

\bibitem{paper:DP-Dixmier}
Marius Dadarlat and Ulrich Pennig.
\newblock A {D}ixmier-{D}ouady theory for strongly self-absorbing
  {$C^*$}-algebras.
\newblock {\em J. Reine Angew. Math.}, 718:153--181, 2016.

\bibitem{paper:FuncEquivKSpec}
Ivo Dell'Ambrogio, Heath Emerson, Tamaz Kandelaki, and Ralf Meyer.
\newblock A functorial equivariant k-theory spectrum and an equivariant
  lefschetz formula, 2011.

\bibitem{paper:DonovanKaroubi}
P.~Donovan and M.~Karoubi.
\newblock Graded {B}rauer groups and {$K$}-theory with local coefficients.
\newblock {\em Inst. Hautes \'{E}tudes Sci. Publ. Math.}, (38):5--25, 1970.

\bibitem{paper:Elmendorf}
A.~D. Elmendorf.
\newblock Systems of fixed point sets.
\newblock {\em Trans. Amer. Math. Soc.}, 277(1):275--284, 1983.

\bibitem{paper:EP-Circle}
David~E. Evans and Ulrich Pennig.
\newblock Equivariant higher {D}ixmier-{D}ouady theory for circle actions on
  {UHF}-algebras.
\newblock {\em Adv. Math.}, 410:Paper No. 108745, 40, 2022.

\bibitem{paper:EvansPennigTwists}
David~E. Evans and Ulrich Pennig.
\newblock Equivariant higher twisted {$K$}-theory of {$SU(n)$} for exponential
  functor twists.
\newblock {\em J. Topol.}, 15(2):896--949, 2022.

\bibitem{paper:FHT3}
Daniel~S. Freed, Michael~J. Hopkins, and Constantin Teleman.
\newblock Loop groups and twisted {$K$}-theory {III}.
\newblock {\em Ann. of Math. (2)}, 174(2):947--1007, 2011.

\bibitem{paper:GabeSzabo}
James Gabe and Gábor Szabó.
\newblock The dynamical {K}irchberg-{P}hillips theorem, 2023.

\bibitem{paper:Grant}
Mark Grant.
\newblock Symmetrized topological complexity.
\newblock {\em J. Topol. Anal.}, 11(2):387--403, 2019.

\bibitem{paper:Hausmann}
Markus Hausmann.
\newblock {$G$}-symmetric spectra, semistability and the multiplicative norm.
\newblock {\em J. Pure Appl. Algebra}, 221(10):2582--2632, 2017.

\bibitem{book:HigsonGuentner}
Nigel Higson and Erik Guentner.
\newblock Group {C}{{\(^{\ast}\)}}-algebras and {K}-theory.
\newblock In {\em Noncommutative geometry. Lectures given at the C. I. M. E.
  summer school, Martina Franca, Italy, September 3--9, 2000}, pages 137--251.
  Berlin: Springer, 2004.

\bibitem{book:Kervaire}
Michael~A. Hill, Michael~J. Hopkins, and Douglas~C. Ravenel.
\newblock {\em Equivariant stable homotopy theory and the {K}ervaire invariant
  problem}, volume~40 of {\em New Mathematical Monographs}.
\newblock Cambridge University Press, Cambridge, 2021.

\bibitem{paper:KaroubiSurvey}
Max Karoubi.
\newblock Twisted {$K$}-theory---old and new.
\newblock In {\em {$K$}-theory and noncommutative geometry}, EMS Ser. Congr.
  Rep., pages 117--149. Eur. Math. Soc., Z\"{u}rich, 2008.

\bibitem{paper:KirchbergPhillips}
Eberhard Kirchberg and N.~Christopher Phillips.
\newblock Embedding of exact {$C^*$}-algebras in the {C}untz algebra
  {$\Cuntz{2}$}.
\newblock {\em J. Reine Angew. Math.}, 525:17--53, 2000.

\bibitem{paper:LashofEquivBundles}
R.~K. Lashof.
\newblock Equivariant bundles.
\newblock {\em Illinois J. Math.}, 26(2):257--271, 1982.

\bibitem{paper:Lind}
John~A. Lind.
\newblock Diagram spaces, diagram spectra and spectra of units.
\newblock {\em Algebr. Geom. Topol.}, 13(4):1857--1935, 2013.

\bibitem{paper:LueckUribe}
Wolfgang L\"{u}ck and Bernardo Uribe.
\newblock Equivariant principal bundles and their classifying spaces.
\newblock {\em Algebr. Geom. Topol.}, 14(4):1925--1995, 2014.

\bibitem{book:MayClass}
J.~Peter May.
\newblock Classifying spaces and fibrations.
\newblock {\em Mem. Amer. Math. Soc.}, 1:xiii+98, 1975.

\bibitem{book:May}
J.~Peter May.
\newblock {\em {$E\sb{\infty }$} ring spaces and {$E\sb{\infty }$} ring
  spectra}, volume Vol. 577 of {\em Lecture Notes in Mathematics}.
\newblock Springer-Verlag, Berlin-New York, 1977.
\newblock With contributions by Frank Quinn, Nigel Ray, and J\o rgen Tornehave.

\bibitem{paper:MayMerlingOsorno}
J.~Peter May, Mona Merling, and Angélica~M. Osorno.
\newblock Equivariant infinite loop space theory, the space level story.
\newblock 2022.

\bibitem{paper:MeyerGroups}
Ralf Meyer.
\newblock On the classification of group actions on {$\rm C^*$}-algebras up to
  equivariant {KK}-equivalence.
\newblock {\em Ann. K-Theory}, 6(2):157--238, 2021.

\bibitem{paper:Murayama}
Mitutaka Murayama.
\newblock On {$G$}-{ANR}s and their {$G$}-homotopy types.
\newblock {\em Osaka J. Math.}, 20(3):479--512, 1983.

\bibitem{book:Naber}
Gregory~L. Naber.
\newblock {\em Topology, geometry, and gauge fields}, volume~25 of {\em Texts
  in Applied Mathematics}.
\newblock Springer, New York, second edition, 2011.
\newblock Foundations.

\bibitem{paper:Palais}
Richard~S. Palais.
\newblock The classification of {$G$}-spaces.
\newblock {\em Mem. Amer. Math. Soc.}, 36:iv+72, 1960.

\bibitem{book:RLL}
M.~R\o{}rdam, F.~Larsen, and N.~Laustsen.
\newblock {\em An introduction to {$K$}-theory for {$C^*$}-algebras}, volume~49
  of {\em London Mathematical Society Student Texts}.
\newblock Cambridge University Press, Cambridge, 2000.

\bibitem{book:RordamStormer}
M.~R\o{}rdam and E.~St\o{}rmer.
\newblock {\em Classification of nuclear {$C^*$}-algebras. {E}ntropy in
  operator algebras}, volume 126 of {\em Encyclopaedia of Mathematical
  Sciences}.
\newblock Springer-Verlag, Berlin, 2002.
\newblock Operator Algebras and Non-commutative Geometry, 7.

\bibitem{paper:Rosenberg}
Jonathan Rosenberg.
\newblock Continuous-trace algebras from the bundle theoretic point of view.
\newblock {\em J. Austral. Math. Soc. Ser. A}, 47(3):368--381, 1989.

\bibitem{paper:Santhanam}
Rekha Santhanam.
\newblock Units of equivariant ring spectra.
\newblock {\em Algebr. Geom. Topol.}, 11(3):1361--1403, 2011.

\bibitem{paper:Schick}
Thomas Schick.
\newblock {$L^2$}-index theorems, {$KK$}-theory, and connections.
\newblock {\em New York J. Math.}, 11:387--443, 2005.

\bibitem{paper:Schlichtkrull}
Christian Schlichtkrull.
\newblock Units of ring spectra and their traces in algebraic {$K$}-theory.
\newblock {\em Geom. Topol.}, 8:645--673, 2004.

\bibitem{book:SchwedeGlobal}
Stefan Schwede.
\newblock {\em Global homotopy theory}, volume~34 of {\em New Mathematical
  Monographs}.
\newblock Cambridge University Press, Cambridge, 2018.

\bibitem{paper:SegalCatAndCoh}
Graeme Segal.
\newblock Categories and cohomology theories.
\newblock {\em Topology}, 13:293--312, 1974.

\bibitem{paper:Strom}
Arne Str\o{}m.
\newblock Note on cofibrations.
\newblock {\em Math. Scand.}, 19:11--14, 1966.

\bibitem{paper:Sullivan}
Dennis~P. Sullivan.
\newblock {\em Geometric topology: localization, periodicity and {G}alois
  symmetry}, volume~8 of {\em $K$-Monographs in Mathematics}.
\newblock Springer, Dordrecht, 2005.
\newblock The 1970 MIT notes, Edited and with a preface by Andrew Ranicki.

\bibitem{paper:TWW}
Aaron Tikuisis, Stuart White, and Wilhelm Winter.
\newblock Quasidiagonality of nuclear {$C^\ast$}-algebras.
\newblock {\em Ann. of Math. (2)}, 185(1):229--284, 2017.

\bibitem{paper:tomDieck}
Tammo tom Dieck.
\newblock Faserb\"undel mit {G}ruppenoperation.
\newblock {\em Arch. Math. (Basel)}, 20:136--143, 1969.

\bibitem{paper:TomsWinter}
Andrew~S. Toms and Wilhelm Winter.
\newblock Strongly self-absorbing {$C^*$}-algebras.
\newblock {\em Trans. Amer. Math. Soc.}, 359(8):3999--4029, 2007.

\bibitem{paper:Trout}
Jody Trout.
\newblock On graded {$K$}-theory, elliptic operators and the functional
  calculus.
\newblock {\em Illinois J. Math.}, 44(2):294--309, 2000.

\bibitem{paper:WanerClassFib}
Stefan Waner.
\newblock Equivariant classifying spaces and fibrations.
\newblock {\em Trans. Amer. Math. Soc.}, 258(2):385--405, 1980.

\bibitem{paper:Waner}
Stefan Waner.
\newblock Equivariant homotopy theory and {M}ilnor's theorem.
\newblock {\em Trans. Amer. Math. Soc.}, 258(2):351--368, 1980.

\end{thebibliography}
